\documentclass[12pt,a4paper,oneside]{report}
\usepackage[T1]{fontenc}
\usepackage[utf8]{inputenc}
\usepackage[english]{babel}
\usepackage{lmodern}
\usepackage[left=32mm,right=25mm,top=27mm,bottom=27mm]{geometry}
\usepackage{setspace}
\onehalfspacing
\usepackage{microtype,amsmath,amssymb,bm}
\usepackage{graphicx,booktabs,tabularx,longtable,array}
\usepackage{enumitem,xcolor,fancyhdr,caption}
\usepackage[numbers,sort&compress]{natbib}
\usepackage{xurl,hyperref}
\usepackage[nameinlink,noabbrev]{cleveref}
\usepackage{bookmark,titlesec,placeins}
\newcommand{\ThesisTitle}{Physics-Informed Modeling of Fluid--Particle Transport in Channels}
\newcommand{\ThesisCoverTitle}{Physics-Informed Modeling of\\Fluid--Particle Transport in Channels}
\newcommand{\ThesisSubtitle}{From Dilute Particle Tracking to Gaussian Parcels and Constrictions}
\newcommand{\ShortTitle}{Physics-Informed Particle Transport}
\newcommand{\AuthorName}{Mohammad Reza Sayah}
\newcommand{\UniversityName}{Sharif University of Technology}
\newcommand{\DepartmentName}{Department of Mechanical Engineering}
\newcommand{\DegreeName}{Master of Science}
\newcommand{\SupervisorName}{Mohammad Ali Bijarchi}
\newcommand{\UniversityCity}{Tehran, Iran}

\newcommand{\FlowVelocity}{1.510}
\newcommand{\FlowPressure}{0.863}
\newcommand{\DragError}{0.717}
\newcommand{\PoiseuilleVelocity}{0.217}
\newcommand{\PoiseuillePressure}{0.180}
\newcommand{\FluxDeviation}{2.266}
\newcommand{\MeshDifference}{0.173}
\newcommand{\InterpolationDifference}{0.177}
\newcommand{\SeedVelocityMean}{1.469}
\newcommand{\SeedVelocitySD}{0.057}
\newcommand{\VOneVelocity}{0.066}
\newcommand{\VOneDistance}{0.024}
\newcommand{\VTwoFluid}{0.096}
\newcommand{\VTwoParticle}{0.238}
\newcommand{\TGU}{0.218}
\newcommand{\TGV}{0.187}
\newcommand{\TGP}{0.215}

\expandafter\def\csname Evidencebasic_errors\endcsname{Plane Poiseuille PINN & Velocity & \PoiseuilleVelocity\% \\
Plane Poiseuille PINN & Centered pressure & \PoiseuillePressure\% \\
Plane Poiseuille FEM & Velocity & $8.192\times10^{-10}$ \\
Curved-channel PINN & Velocity & \FlowVelocity\% \\
Curved-channel PINN & Centered pressure & \FlowPressure\% \\
ForceNet & Normalized drag & \DragError\% \\
Section flux & Maximum relative deviation & \FluxDeviation\% \\
Two finest FEM meshes & Velocity difference & \MeshDifference\% \\
}
\expandafter\def\csname Evidencemesh\endcsname{290 & 1298 & 180 & 1.141 & $6.735\times10^{-11}$ \\
1298 & 5474 & 720 & 0.173 & $1.097\times10^{-10}$ \\
3026 & 12530 & 1620 & Reference & $1.600\times10^{-10}$ \\
}
\expandafter\def\csname Evidenceseeds\endcsname{42 & 1.510 & 0.863 \\
43 & 1.429 & 0.886 \\
}
\expandafter\def\csname Evidencegaussian_errors\endcsname{V01 & Stokes velocity & \VOneVelocity\% \\
V01 & Integrated displacement & \VOneDistance\% \\
V02 & Fluid velocity & \VTwoFluid\% \\
V02 & Particle velocity & \VTwoParticle\% \\
V10 & Taylor--Green $u$ & \TGU\% \\
V10 & Taylor--Green $v$ & \TGV\% \\
V10 & Taylor--Green $p$ & \TGP\% \\
}
\expandafter\def\csname Evidencecoupled_scope\endcsname{0.08 & 0.05 & 0.3854 & 0.0476 & 12.85 \\
0.20 & 0.05 & 0.1372 & 0.0267 & 4.57 \\
0.20 & 0.15 & 0.3526 & 0.0834 & 11.75 \\
0.60 & 0.05 & 0.0407 & 0.0017 & 1.36 \\
}
\expandafter\def\csname Evidencetrajectories\endcsname{-0.650 & 2.67 & 98.86 & Yes \\
-0.433 & 2.30 & 99.33 & Yes \\
-0.217 & 4.30 & 99.27 & Yes \\
0.000 & 7.83 & 98.59 & Yes \\
0.217 & 16.09 & 96.19 & Yes \\
0.433 & 34.97 & 87.39 & Yes \\
0.650 & 3.29 & 73.55 & No \\
}
\expandafter\def\csname Evidenceparcel_parameters\endcsname{$\Rey_H,\ R,\ b$ & $20,\ 1000,\ 1$ \\
$\beta$, parcel count, seed & $0.05,\ 6,\ 7$ \\
$\sigma,\ T$ & $0.14,\ 1$ \\
$\vect{g}_{\rm eff}$ & $(0,\ -0.1)$ \\
$\St,\ d/H$ & $0.60,\ 0.023238$ \\
Multiplicity $N_i$ & $3.804966$ \\
Represented mass $m_i$ & $0.025000$ \\
Single-sphere mass in the planar scale & $0.00657036$ \\
Initial parcel velocity & $(0,0)$ \\
}
\expandafter\def\csname Evidenceparcel_initial_centers\endcsname{P1 & 0.887605 & 0.698607 \\
P2 & 1.053254 & 0.362604 \\
P3 & 0.530183 & 0.686777 \\
P4 & 0.205792 & 0.660614 \\
P5 & 1.076776 & 0.483967 \\
P6 & 0.533336 & 0.389213 \\
}

\newcommand{\EvidenceRows}[1]{\csname Evidence#1\endcsname}
\definecolor{ThesisBlue}{HTML}{173F5F}
\definecolor{ThesisTeal}{HTML}{087F8C}
\definecolor{ThesisGrey}{HTML}{738390}
\hypersetup{colorlinks=true,linkcolor=ThesisBlue,citecolor=ThesisBlue,urlcolor=ThesisBlue,
 pdftitle={\ThesisTitle: \ThesisSubtitle},pdfauthor={\AuthorName},
 pdfsubject={Master of Science thesis manuscript on physics-informed particle transport}}
\setcounter{secnumdepth}{2}\setcounter{tocdepth}{1}
\setlength{\headheight}{16pt}
\pagestyle{fancy}\fancyhf{}
\fancyhead[L]{\small\nouppercase{\leftmark}}\fancyfoot[C]{\thepage}
\renewcommand{\headrulewidth}{0.3pt}
\fancypagestyle{plain}{\fancyhf{}\fancyfoot[C]{\thepage}\renewcommand{\headrulewidth}{0pt}}
\renewcommand{\chaptermark}[1]{\markboth{\thechapter.\ #1}{}}
\titleformat{\chapter}[display]{\normalfont\raggedright\color{ThesisBlue}}{\Large\chaptertitlename\ \thechapter}{8pt}{\huge\bfseries}
\titlespacing*{\chapter}{0pt}{-10pt}{22pt}
\titleformat{\section}{\large\bfseries\color{ThesisBlue}}{\thesection}{.7em}{}
\titleformat{\subsection}{\normalsize\bfseries\color{ThesisTeal}}{\thesubsection}{.7em}{}
\captionsetup{font=small,labelfont=bf,labelsep=period}
\renewcommand{\arraystretch}{1.18}
\setlist[itemize]{topsep=4pt,itemsep=2pt,parsep=0pt}
\setlist[enumerate]{topsep=4pt,itemsep=3pt,parsep=0pt}
\newcolumntype{Y}{>{\raggedright\arraybackslash}X}
\newcommand{\vect}[1]{\bm{#1}}
\newcommand{\Rey}{\mathrm{Re}}
\newcommand{\St}{\mathrm{St}}
\newcommand{\MSE}{\operatorname{MSE}}
\newcommand{\FrontChapter}[1]{\chapter*{#1}\phantomsection\addcontentsline{toc}{chapter}{#1}\markboth{#1}{#1}}
\newcommand{\FigureShort}[1]{\csname FigureShort#1\endcsname}
\expandafter\def\csname FigureShort01_study_route\endcsname{Study progression and proposed constriction}
\expandafter\def\csname FigureShort02_baseline_architecture\endcsname{Baseline flow and force network architecture}
\expandafter\def\csname FigureShort03_poiseuille\endcsname{Analytical plane Poiseuille verification}
\expandafter\def\csname FigureShort04_curved_fields\endcsname{Curved-channel fields and discrepancies}
\expandafter\def\csname FigureShort05_heldout_agreement\endcsname{Independent field and normalized drag comparisons}
\expandafter\def\csname FigureShort06_reference_and_seeds\endcsname{Reference-mesh and training-seed sensitivity}
\expandafter\def\csname FigureShort07_flux\endcsname{Section flux in the curved baseline}
\expandafter\def\csname FigureShort08_gaussian_architecture\endcsname{Coupled fluid and Gaussian parcel representation}
\expandafter\def\csname FigureShort09_stokes_response\endcsname{V01: reduced Stokes response}
\expandafter\def\csname FigureShort10_two_way_response\endcsname{V02: uniform two-way relaxation}
\expandafter\def\csname FigureShort11_taylor_green\endcsname{V10: Taylor--Green carrier-flow verification}
\expandafter\def\csname FigureShort12_kernel_verification\endcsname{Gaussian overlap and bounded-domain normalization}
\expandafter\def\csname FigureShort13_constriction_geometry\endcsname{Proposed constriction geometry and height profiles}
\expandafter\def\csname FigureShort14_boundary_kernel\endcsname{Illustration of kernel support near the throat wall}
\expandafter\def\csname FigureShort15_parcel_transport\endcsname{Computed fluid--parcel transport snapshots}
\expandafter\def\csname FigureShort16_parcel_paths_velocities\endcsname{Parcel trajectories and velocity histories}
\expandafter\def\csname FigureShort17_gaussian_clouds\endcsname{Evolution of represented Gaussian clouds}
\expandafter\def\csname FigureShort18_parcel_carrier_response\endcsname{Computed carrier response at the final time}
\expandafter\def\csname FigureShort19_parcel_balances\endcsname{Energy, drag dissipation and balance discrepancies}
\newcommand{\ResultFigure}[4][.98]{\begin{figure}[!htbp]\centering
 \includegraphics[width=#1\textwidth,height=.67\textheight,keepaspectratio]{figures/#2.pdf}
 \caption[\FigureShort{#2}]{#3}\label{#4}\end{figure}}
\setlength{\emergencystretch}{2em}
\setlength{\parskip}{2pt}
\begin{document}
\hypersetup{pageanchor=false}
\begin{titlepage}
\centering
\vspace*{3mm}
{\Large\UniversityName\par}
\vspace{3mm}{\large\DepartmentName\par}
\vspace{19mm}
{\Huge\bfseries\color{ThesisBlue}\ThesisTitle\par}
\vspace{7mm}
{\Large\color{ThesisTeal}\ThesisSubtitle\par}
\vspace{14mm}
{\Large\AuthorName\par}
\vspace{12mm}
{\large Master's Thesis\par}
\vspace{3mm}{\large\DegreeName\par}
\vspace{12mm}
{\large Supervisor: \SupervisorName\par}
\vfill
{\normalsize\UniversityCity\par\SubmissionDate\par}
\end{titlepage}
\pagenumbering{roman}\hypersetup{pageanchor=true}
\FrontChapter{Abstract}
Predicting particle transport requires an accurate carrier flow, a consistent
interphase force model, and an evaluation that distinguishes field agreement
from reliable particle motion. I develop a staged physics-informed framework
that begins with dilute particle tracking in a curved channel, introduces
Gaussian parcel coupling in a straight channel, and formulates an extension
to a smoothly constricted channel. The baseline combines a continuous
velocity--pressure network with a neural approximation of the prescribed
Stokes-drag law. Independent finite-element and analytical comparisons give
sample-relative errors of \FlowVelocity\% in curved-channel velocity and
\FlowPressure\% in gauge-aligned pressure. The normalized drag error is
\DragError\%. The Gaussian formulation uses the same normalized kernel for
fluid interpolation and reaction-force spreading, yielding exact discrete
force and work identities on the transfer quadrature. Reduced-model tests
recover Stokes relaxation, two-way momentum exchange, Gaussian overlap, and
Taylor--Green decay. For these tests, the reported relative errors are below
0.25\%; this bound concerns the listed solution observables, not every
residual or every channel prediction. Four coupled Gaussian channel
approximations are computed and evaluated separately. The selected case
has a momentum RMS of $0.0407$ against a $0.03$ screen, while its parcel
kinematic and acceleration RMS values are below $0.002$. Baseline trajectory
accuracy varies with initial condition: three of seven paths meet all
position, time-coverage and endpoint criteria. The final stage develops a
proposed constriction study with consistent wall treatment and a
reference-based evaluation protocol. The central contribution is a verified foundation for
progressing from one-way particle tracking to conservative regularized
coupling, together with explicit conditions for extending that foundation
to confined geometries.

\vspace{5mm}\noindent\textbf{Keywords:} physics-informed neural networks;
particle transport; Stokes drag; Gaussian parcels; two-way coupling;
channel constriction; numerical verification.

\clearpage\phantomsection\addcontentsline{toc}{chapter}{Contents}\tableofcontents
\clearpage\phantomsection\addcontentsline{toc}{chapter}{List of Figures}\listoffigures
\clearpage\phantomsection\addcontentsline{toc}{chapter}{List of Tables}\listoftables
\FrontChapter{Abbreviations and Nomenclature}
\begingroup\small\setstretch{1.02}
\begin{tabularx}{\textwidth}{@{}lY@{}}\toprule
Abbreviation & Meaning\\\midrule
AD & Automatic differentiation\\
CFD & Computational fluid dynamics\\
ERPP & Exact regularized point-particle method\\
FEM & Finite-element method\\
MLP & Multilayer perceptron\\
ODE / PDE & Ordinary / partial differential equation\\
PINN & Physics-informed neural network\\
RMS & Root-mean-square\\
\bottomrule\end{tabularx}
\vspace{7mm}
\begin{longtable}{@{}p{.22\textwidth}p{.73\textwidth}@{}}
\toprule Symbol & Meaning\\\midrule\endhead
$\vect{u},p$ & Dimensionless fluid velocity and pressure\\
$\vect{q}_i,\vect{v}_i$ & Parcel center and parcel velocity\\
$\vect{U}_i$ & Fluid velocity interpolated to parcel $i$\\
$\vect{F}_i,\vect{S}_a$ & Represented hydrodynamic drag and reaction source\\
$L_0,U_0,H,B$ & Reference length, speed, channel height and spanwise depth\\
$b=B/H$ & Reference-depth ratio for planar sphere parcels\\
$d_p,\delta_p$ & Particle diameter scaled by $L_0$ or $H$, respectively\\
$R=\rho_p/\rho_f$ & Particle-to-fluid density ratio\\
$\Rey,\Rey_p,\St$ & Flow Reynolds, slip Reynolds and Stokes numbers\\
$\beta,\alpha$ & Mass loading and effective local solid volume fraction\\
$N_i,m_i$ & Parcel multiplicity and represented dimensionless mass\\
$\sigma,W_{ai},Z_i$ & Gaussian width, normalized transfer weight and normalizer\\
$\Delta A_a$ & Transfer-quadrature area associated with point $a$\\
$s,\eta$ & Centerline coordinate and normalized transverse coordinate\\
$Q,h,H_t$ & Volume flux per reference depth, half-height and throat height\\
$\delta,\ell,x_c$ & Constriction severity, half-length and center\\
$E_u,E_p$ & Discrete relative velocity and gauge-aligned pressure errors\\
\bottomrule
\end{longtable}
\noindent Pressure uses a viscous scale in the baseline and an inertial scale
in the Gaussian channel. Each formulation states its scale explicitly.
The Gaussian width $\sigma$ is distinct from both the particle radius and
the channel half-height $h$.
\endgroup
\clearpage\pagenumbering{arabic}
\chapter{Introduction}\label{ch:introduction}

\section{Motivation}
Particle-laden channel flows connect local hydrodynamics with transport
outcomes. A particle can follow the carrier flow closely, cross a streamline
because of inertia, or approach a wall closely enough to change its final
destination. In a narrow passage, these outcomes depend on both the flow
field and the force closure. A small average velocity error can coexist
with a significant error in a particle's residence time or contact location.
Predictive modeling therefore requires more than visually plausible
streamlines.

Conventional computational methods provide established ways to discretize
the incompressible Navier--Stokes equations and integrate particle motion
\citep{ferziger2002,clift1978}. Physics-informed neural networks offer a
different representation: a differentiable network approximates the
solution, while differential equations and boundary conditions constrain
its parameters \citep{raissi2019,karniadakis2021}. This representation is
attractive when smooth field evaluation and derivatives are needed at
many moving particle positions. Its utility must nevertheless be established
by comparison with an independent reference and by checks that are sensitive
to the intended transport quantities.

I organize this study around a progressive increase in physical coupling
and geometric confinement. The first stage considers dilute particles whose
motion does not alter the carrier flow. The second introduces represented
particle populations and a spatially regularized reaction force in a straight
channel. The third defines a smooth constriction as the next geometry to be
studied. This progression isolates the consequences of adding two-way
coupling before introducing a varying channel height.

\section{Problem statement}
The central problem is to determine which parts of a physics-informed
fluid--particle model have sufficient numerical support for transport
prediction. There are three distinct questions. The first concerns the
carrier field: does the predicted velocity and pressure agree with a solution
of the same boundary-value problem? The second concerns the force transfer:
does the discrete coupling apply equal and opposite hydrodynamic forces
without introducing artificial interphase power? The third concerns
particle transport: are trajectories, contact events, and transit times
accurate for the geometry under consideration?

These questions are related but are not interchangeable. An algebraically
conservative transfer operator does not eliminate the residual error of a
trained fluid field. A low drag-network error does not correct an inaccurate
near-wall velocity. A successful straight-channel benchmark does not
establish performance at a constriction. I therefore treat each claim at
the level at which it has been evaluated.

\section{Aim and objectives}
The aim is to establish a physically consistent and numerically assessable
framework for physics-informed particle transport in channels. The study
has four objectives:
\begin{enumerate}
\item Formulate and evaluate a one-way coupled baseline with independent
carrier-flow and analytical drag references.
\item Derive a Gaussian parcel coupling that preserves force and interphase
work on its transfer quadrature.
\item Evaluate the components of the coupled formulation using independent
analytical and numerical verification problems.
\item Specify a constricted-channel extension whose geometry, boundary
conditions, and acceptance criteria can be tested without altering the
meaning of the earlier benchmarks.
\end{enumerate}

The corresponding research questions are: How accurately does the baseline
recover independently computed channel flow? Which conservation properties
follow from the Gaussian discretization, and which require additional
solution accuracy? What evidence is still needed before the formulation
can support predictions in a constricted channel?

\section{Scope and physical assumptions}
The completed channel models describe incompressible Newtonian flow and
spherical heavy particles moving in a plane. The baseline uses one-way
Stokes drag and a prescribed irreversible wall-sticking rule. The Gaussian
model represents spherical-particle populations through regularized force
distributions. The Gaussian cloud is a force-transfer representation; it
is not a resolved solid sphere and it is not an independently validated
description of a dense particle agglomerate.

The two-dimensional carrier does not resolve spanwise motion. No inference
about a finite-depth three-dimensional channel follows from an extruded
view. Particle rotation, lubrication, history force, and hydrodynamic
interactions of resolved finite-size bodies lie outside the accepted
baseline. The constriction is considered as a static wall geometry rather
than a dynamically growing deposit or a biological vessel.

\section{Study sequence and contributions}
\Cref{fig:route} summarizes the study sequence. The first contribution is a
baseline evaluation that combines analytical plane flow, an independently
assembled finite-element reference, a held-out drag test, and an explicit
distinction between field and trajectory accuracy. The second is a
consistent planar sphere-parcel formulation with shared Gaussian
interpolation and spreading. The third is a verification chain connecting
reduced relaxation problems, normalization identities, and an unsteady
incompressible-flow benchmark. The final contribution is a defined
constricted-channel study design that preserves the same requirements for
reference matching and conservation.

Completed findings are presented for the baseline carrier field and the
Gaussian model's verification problems. The coupled Gaussian channel
predictions have unresolved residual errors. The constricted-channel
formulation is prospective. This division determines the interpretation of
the conclusions and prevents a proposed geometry from being presented as
an executed calculation.

\ResultFigure[.92]{01_study_route}{Study progression. The first two panels
identify the model classes investigated; the third defines the proposed
constricted-channel extension. All three panels are geometric illustrations,
not simulation fields. Gaussian shading depicts kernel support rather than
a measured particle concentration.}{fig:route}

\section{Thesis organization}
\Cref{ch:literature} reviews the modeling choices and the requirements for
credible numerical evidence. \Cref{ch:framework} defines the common
nondimensional and evaluation framework. \Cref{ch:basic} develops the
baseline particle problem and reports its supported results.
\Cref{ch:gaussian} introduces the straight-channel Gaussian formulation and
its verification results. \Cref{ch:constriction} defines the constricted
channel and the proposed evaluation protocol. \Cref{ch:discussion} compares
the implications of the stages, and \cref{ch:conclusions} gives the
conclusions and priorities for further work. Supplementary derivations and
the computational record appear in the appendices.

\chapter{Literature Review}\label{ch:literature}

\section{Modeling levels for particle-laden flow}
The hydrodynamics of a particle suspension can be described at different
levels of resolution. A particle-resolved method imposes boundary conditions
on each particle surface. It can represent the disturbance flow, finite-size
effects, and surface stresses when these are sufficiently resolved. Direct
forcing immersed-boundary methods provide one established route to such
calculations \citep{uhlmann2005}. Their additional resolution requirements
are justified when particle-scale flow is itself an essential part of the
problem.

A point-particle model instead advances particle centers using a force law
obtained under specified physical assumptions. The classical equation of
motion for a small rigid sphere includes effects beyond steady drag
\citep{maxey1983}. The heavy-particle, low-slip-Reynolds-number approximation
used here retains a much narrower description. It is appropriate to assess
this approximation as a simplified model rather than attribute to it the
full physics of the parent equation.

The strength of a point-particle approach is its economy of representation.
Its limits follow from the same reduction. A particle center and diameter
do not impose no-slip conditions on an actual moving surface. Finite-size
wall stresses, rotational dynamics, and lubrication do not appear unless
additional closures are included. These distinctions are particularly
important when a channel becomes narrow relative to particle diameter
\citep{clift1978}.

\section{One-way and two-way coupling}
In one-way coupling, the fluid field determines the particle force but the
particle reaction is neglected in the fluid equations. This separation is
useful for studying dilute transport because the fluid solution can be
evaluated independently of particle initial conditions. Small particle
volume fraction is a necessary consideration, although it does not alone
guarantee negligible feedback when the particle density is high. Mass
loading and the relevant time scales also matter.

Two-way coupling includes the hydrodynamic reaction in the carrier-fluid
momentum balance. The coupling must specify how a particle samples the fluid
and how its reaction is distributed to the Eulerian field. Treating either
operation independently can generate an inconsistent power exchange even
when the integrated force is correct. A shared transfer operator provides a
direct means to avoid this particular inconsistency.

Regularized point-particle methods provide a principled context for this
transfer. In the exact regularized point-particle method, the treatment of
the disturbance and its regularization is tied to the governing flow
equations \citep{gualtieri2015}. A generic normalized Gaussian deposition
does not inherit all of that method's properties. In particular, sampling
the disturbed velocity at a force source is not equivalent to recovering
the undisturbed velocity required by an isolated-particle drag law.

\section{Gaussian parcels and represented populations}
A parcel represents a population of physical particles through a
multiplicity. Its represented mass must scale with that multiplicity, and
its reaction force must use the same convention. A Gaussian kernel spreads
that force over a finite region, avoiding a singular source in the carrier
equations. The width of the kernel is a modeling parameter with a direct
effect on the spatial distribution of forcing.

Three lengths must be kept distinct: the physical diameter, the numerical
or model regularization width, and the local channel height. Equating the
Gaussian width with a particle radius without justification would blur
point-particle and finite-size interpretations. Likewise, changing the
width while keeping quadrature resolution fixed does not constitute an
ordinary grid-convergence study: it changes the regularized model being
solved.

Near a wall, a free-space Gaussian is truncated by the fluid domain. A
discrete normalization can preserve its integrated weight on the chosen
quadrature. That normalization establishes a transfer identity; it does
not derive a wall-corrected drag law. This distinction carries directly into
the constricted-channel extension developed later.

\section{Physics-informed representations}
Physics-informed neural networks approximate solution fields with
differentiable functions and train those functions using equation and
boundary residuals \citep{raissi2019}. They provide continuous field values
without a separate spatial interpolant. Automatic differentiation also
allows the same representation to supply the derivatives entering the
equations. These features are useful for particle transport, where the
evaluation locations move continuously.

Physics information in the objective is not a guarantee that optimization
finds a sufficiently accurate solution. Residual terms can have different
scales and different optimization behavior. Gradient imbalance and training
difficulties are documented limitations of PINNs
\citep{wang2021,wang2022}. A weighted objective may decrease while a
transport-sensitive region remains inaccurate. Consequently, reporting only
the optimization history would provide insufficient support for the
predictions considered here.

Hard constraints offer an additional design choice. An ansatz can impose an
initial condition, periodicity, or homogeneous wall velocity exactly, leaving
the network to represent a correction that respects those conditions. This
reduces the number of requirements enforced through a penalty. It does not
remove the need to evaluate momentum, continuity, and particle-motion
residuals at independent points.

The baseline and Gaussian stages use different neural roles. The baseline
flow network approximates a steady velocity--pressure field, whereas the
force network approximates a prescribed algebraic law. In the Gaussian
stage, separate fluid and parcel networks are coupled through common
residuals. This is a continuous-time coupled forward representation rather
than a trajectory regression against reference data.

\section{Independent references and numerical credibility}
The finite-element method supplies an independently discretized reference
for the baseline carrier flow. Taylor--Hood velocity--pressure elements are
a standard mixed formulation for incompressible flow \citep{taylor1973}.
Their assembly here uses \texttt{scikit-fem} \citep{gustafsson2020}. An
independent discretization is valuable because its numerical errors differ
from those of the neural representation. It must still be evaluated for
nonlinear convergence, incompressibility, and mesh sensitivity.

Verification and physical validation address different questions
\citep{roache1998}. Agreement with an analytical solution tests whether a
specified mathematical problem is solved correctly. Agreement with another
numerical solver tests consistency for a matched model and boundary-value
problem. Experimental validation additionally requires suitable
measurements, controlled conditions, and compatible physical assumptions.
The results in this thesis primarily concern the first two questions.

A relative velocity error should be interpreted together with local
discrepancies, boundary errors, flux variation, and particle-motion
diagnostics. Pressure requires particular care because incompressible
velocity boundary conditions determine pressure only up to an additive
constant. Comparing unaligned pressure arrays can produce a large error
without a corresponding physical disagreement. Conversely, an error
measured on a section where the reference pressure is constant can have an
ill-conditioned denominator. A spatially extended, gauge-aligned comparison
avoids these problems.

\section{Transport accuracy and confinement}
Particle endpoints depend on a time-integrated response to the carrier
field. A small discrepancy may accumulate along a long path or determine
whether a particle approaches a capture boundary. Near-wall outcomes can
therefore be less accurate than bulk flow metrics suggest. A meaningful
trajectory comparison must pair the same initial conditions, restrict its
position error to the common time interval, and separately report whether
the two solutions reach the same endpoint class.

Confinement adds a physical limitation to this numerical sensitivity. When
the local channel height approaches the particle diameter, the unbounded
Stokes law becomes a questionable description even if the numerical
trajectory is well resolved. A study of narrowing should therefore track
both numerical agreement and the changing ratio of particle diameter to
local gap. A dilute regularized model cannot be used to infer dense
blockage merely because a narrow geometry is drawn.

\section{Position of the present study}
\Cref{tab:model_comparison} summarizes the modeling levels relevant to the
study. The purpose is to build a verified chain across the first two levels
and specify the requirements for a constrained extension. The focus is
neither a new empirical drag correlation nor a replacement for resolved
particle hydrodynamics. The main methodological question is whether a
neural representation can support the stated fluid and particle quantities
while preserving a consistent transfer of force and work.

\begin{table}[htbp]\centering\small
\caption{Model classes and the evidence they require.}\label{tab:model_comparison}
\begin{tabularx}{\textwidth}{@{}p{.22\textwidth}YY@{}}\toprule
Model & Main representation & Essential assessment\\\midrule
One-way point particles & Carrier field and particle ODEs & Field error, slip regime, trajectory agreement\\
Gaussian parcels & Population-weighted regularized forcing & Shared transfer, residuals, budgets, width sensitivity\\
Resolved particles & Moving solid surfaces and surrounding flow & Surface resolution, forces, torques and gap flow\\
Constricted dilute model & Variable-height fluid domain with particles & Wall geometry, flux, throat resolution and closure validity\\
\bottomrule\end{tabularx}
\end{table}

\chapter{Mathematical and Evaluation Framework}\label{ch:framework}

\section{Baseline nondimensionalization}
Let $L_0$ and $U_0$ denote reference length and velocity, and let $\rho_f$
and $\mu$ denote the fluid density and dynamic viscosity. The baseline
uses
\begin{equation}
 \vect{x}^{*}=\frac{\vect{x}}{L_0},\quad
 t^{*}=\frac{tU_0}{L_0},\quad
 \vect{u}^{*}=\frac{\vect{u}}{U_0},\quad
 p_v^{*}=\frac{pL_0}{\mu U_0},\quad
 \Rey=\frac{\rho_fU_0L_0}{\mu}.
 \label{eq:baseline_scales}
\end{equation}
Stars are suppressed after the scales have been specified. The steady
carrier equations are
\begin{equation}
 \nabla\cdot\vect{u}=0,\qquad
 \Rey(\vect{u}\cdot\nabla)\vect{u}
 =-\nabla p_v+\nabla^2\vect{u}.
 \label{eq:baseline_ns}
\end{equation}
Using a viscous pressure scale places $\Rey$ on the convective term.
Pressure values from this model cannot be compared directly with a model
using an inertial pressure scale without conversion.

For a spherical heavy particle, the retained force is the creeping-flow
drag. With $R=\rho_p/\rho_f$ and $d_p=d/L_0$, its dimensionless equations
are
\begin{equation}
 \frac{d\vect{q}}{dt}=\vect{v},\qquad
 \frac{d\vect{v}}{dt}
 =\frac{18}{R\Rey d_p^2}[\vect{u}(\vect{q})-\vect{v}]
 =\frac{\vect{u}(\vect{q})-\vect{v}}{\St},
 \qquad \St=\frac{R\Rey d_p^2}{18}.
 \label{eq:baseline_particle}
\end{equation}
The slip Reynolds number is
\begin{equation}
 \Rey_p=\Rey d_p|\vect{u}(\vect{q})-\vect{v}|.
 \label{eq:slip_re}
\end{equation}
The physical rationale for Stokes drag requires $\Rey_p\ll1$, along with a
small particle relative to the spatial scales of the carrier flow.

\section{Geometry and boundary-value matching}
The baseline channel is a constant-width tube around a planar centerline
$\vect{c}(s)$. Let $\vect{n}(s)$ denote its unit normal and let $h$ denote
the half-width. Coordinates are mapped by
\begin{equation}
 \vect{x}(s,\eta)=\vect{c}(s)+h\eta\vect{n}(s),\qquad
 0\leq s\leq L,\quad -1\leq\eta\leq1.
 \label{eq:centerline_map}
\end{equation}
The carrier velocity vanishes on both walls. At each end section, a
parabolic velocity in the local tangent direction is prescribed:
\begin{equation}
 \vect{u}(s_{\mathrm{end}},\eta)
 =(1-\eta^2)\vect{t}(s_{\mathrm{end}}).
 \label{eq:end_bc}
\end{equation}
Thus the reference problem uses velocity conditions at both ends, rather
than an imposed pressure outlet. A pressure gauge fixes the remaining
constant. The neural and finite-element comparisons must use this same
problem.

The associated volume flux per unit reference depth is
\begin{equation}
 Q_0=h\int_{-1}^{1}(1-\eta^2)\,d\eta=\frac{4h}{3}.
 \label{eq:flux}
\end{equation}
For a steady incompressible flow without sources of fluid mass, the section
flux should be independent of $s$. This makes flux variation a useful
integrated diagnostic alongside local differential residuals.

\section{Gaussian-channel scales}
For the Gaussian stage, the reference length is the channel height $H$.
The pressure scale is inertial:
\begin{equation}
 p_i^{*}=\frac{p}{\rho_f U_0^2},\qquad
 \Rey_H=\frac{U_0H}{\nu},\qquad \delta_p=\frac{d}{H}.
 \label{eq:gaussian_scales}
\end{equation}
For identical length and velocity scales, the two pressure conventions
satisfy $p_v^{*}=\Rey_H p_i^{*}$. The governing equation is correspondingly
written with $1/\Rey_H$ multiplying viscous diffusion.

The carrier is two-dimensional, but the represented physical particles are
spheres. A reference depth $B$ is therefore required to convert represented
sphere mass into a planar source. With $b=B/H$ and multiplicity $N_i$, the
dimensionless represented mass is
\begin{equation}
 m_i=\frac{N_iR\pi\delta_{p,i}^3}{6b},\qquad
 \St_i=\frac{R\Rey_H\delta_{p,i}^2}{18},\qquad
 \beta=\frac{\sum_i m_i}{|\Omega|}.
 \label{eq:parcel_mass}
\end{equation}
This convention scales physical mass by $\rho_fH^2B$. The kernel has units
of inverse area before nondimensionalization. Replacing it with a
three-dimensional volume kernel without changing the mass convention would
be dimensionally inconsistent.

\section{What is learned}
The neural representations approximate fields or trajectories; the
physical equations specify their constraints. The baseline force network
learns a known drag response from an algebraic residual. It does not infer
an unknown force law from experiments. The Gaussian parcel network predicts
time-dependent position and velocity corrections whose consistency is
enforced by both kinematic and dynamic residuals.

Adam is used for initial optimization, with optional limited-memory
quasi-Newton refinement \citep{kingma2015,liu1989}. The optimizer sequence
does not define the credibility of the final state. Checkpoint selection
uses a fixed residual monitor, whereas the reported field comparison uses
an independent reference. A longer training budget is a computational
choice, not an accuracy certificate.

\section{Error measures and pressure gauge}
For values at $M$ independent locations, the sample-relative velocity error
is
\begin{equation}
 E_u=\left[
 \frac{\sum_{j=1}^{M}|\vect{u}^{\mathrm{PINN}}_j-
               \vect{u}^{\mathrm{ref}}_j|^2}
      {\sum_{j=1}^{M}|\vect{u}^{\mathrm{ref}}_j|^2}
 \right]^{1/2}.
 \label{eq:velocity_error}
\end{equation}
These are discrete norms over the specified samples. Unless area weights
are explicitly included, they are not exact domain-integral norms. For
pressure, I first remove the arithmetic mean of each compared array:
\begin{equation}
 \widetilde p_j=p_j-\frac{1}{M}\sum_{k=1}^{M}p_k,\qquad
 E_p=\frac{\|\widetilde{\vect{p}}^{\mathrm{PINN}}-
                    \widetilde{\vect{p}}^{\mathrm{ref}}\|_2}
                   {\|\widetilde{\vect{p}}^{\mathrm{ref}}\|_2}.
 \label{eq:pressure_error}
\end{equation}
The pressure samples span the streamwise direction so that the reference
denominator is meaningful. Mean subtraction removes only a constant; it
does not rescale a pressure gradient or modify the predicted dynamics.

The section-flux diagnostic is
\begin{equation}
 E_Q=\max_s\left|\frac{Q(s)}{Q_0}-1\right|.
 \label{eq:flux_error}
\end{equation}
Reported differential residuals use their stated nondimensional equations.
Their magnitudes should not be compared across pressure scales, equation
normalizations, or geometries without accounting for those differences.

\section{Trajectory evaluation}
Consider a neural and reference trajectory with the same initial condition.
If their common interval is $[0,T_c]$, the position discrepancy is measured
over that interval and normalized by channel width $W=2h$. Endpoint
agreement is assessed separately using the terminal event class. I also
retain the fraction of the longer trajectory covered by the common
interval. A low error calculated over a short common interval is not
sufficient evidence of agreement over the full transit.

Integration success has a narrower meaning: it indicates that the numerical
ODE solver completed its assigned calculation without reporting failure.
It does not establish that the neural velocity or force was accurate.
Accordingly, the acceptance criteria distinguish integration status,
position discrepancy, common-time coverage, and endpoint agreement.

\section{Evidence hierarchy}
I use four levels of numerical evidence. Algebraic identities establish
properties of the coupling operators. Analytical benchmarks test reduced
equations with known solutions. Independently discretized references test
the field approximation for a matched problem. Finally, transport
comparisons assess the accumulated consequences of the field and force
errors. None of these alone establishes experimental agreement.

For the baseline, the numerical screening limits are 5\% for field and
normalized drag errors, 2\% for boundary RMS and reference-interpolation
error, and 5\% for maximum flux deviation. Baseline particle acceptance is
stricter in its combination of conditions: every tested endpoint class must
agree, every width-normalized RMS position error must be below 5\%, and
every common-time fraction must exceed 95\%. For the Gaussian verification
tests, the stated relative-error and independent residual criteria are
0.03. These limits are study criteria rather than universal physical
thresholds. Their actual outcomes are reported where the corresponding
results are presented.

\chapter{Baseline Particle Problem}\label{ch:basic}

\section{Physical problem and parameter set}
The baseline separates carrier-flow approximation from particle response.
The carrier is steady and two-dimensional. Particles are spherical, heavy,
and dilute, and their forces do not feed back into the fluid. The principal
channel consists of two connected semicircular bends with nominal inner
and outer radii $R_{\mathrm{in}}=0.2$ and $R_{\mathrm{out}}=0.5$. Its
width is $W=0.3$ and its centerline length is approximately $2.1991$, all
scaled by $L_0$. The Reynolds number is $\Rey=1$.

The baseline particle parameters are $R=1000$ and $d_p=0.025$. The resulting
Stokes number is $\St\simeq0.03472$. The channel wall rule stops the
particle when its center approaches within one radius of a wall. This is
a specified irreversible sticking model, not a prediction of adhesion
chemistry or a resolved impact process. In the full trajectory set, the
largest recorded slip Reynolds number is approximately $0.01221$, which
supports the low-slip-Reynolds-number assumption for these calculations.
The ratio $d_p/W\simeq0.0833$ still warrants caution for near-wall motion.

\section{Flow and force representations}
\subsection{FlowNet}
The flow network maps five input slots to velocity and pressure:
\begin{equation}
 (x,y,R_{\mathrm{in}},R_{\mathrm{out}},\Rey)
 \longmapsto (u,v,p_v).
 \label{eq:flownet}
\end{equation}
Velocity--pressure PINN formulations for incompressible flow provide
related methodological background \citep{jin2021}. The present network
has five hidden layers with 70 tanh units per layer. % Citation-plan C19
 The spatial inputs
are centered and scaled, and the pressure output is scaled with a
channel-based pressure magnitude. The geometry and Reynolds inputs are
fixed within an individual trained model. Their presence does not establish
generalization across untrained geometries or Reynolds numbers.

The objective combines momentum and continuity residuals, wall and end
velocity conditions, a pressure gauge, and a section-flux condition. If
$r_c=\nabla\cdot\vect{u}$ and $\vect{r}_m$ denotes the residual of
\cref{eq:baseline_ns}, the scaled physics term is
\begin{equation}
 \mathcal{L}_{\mathrm{phys}}
 =\left\langle (hr_c)^2\right\rangle
 +\left\langle |h^2\vect{r}_m|^2\right\rangle.
 \label{eq:baseline_loss}
\end{equation}
Here the averages follow the componentwise mean-square convention of the
implementation. The complete objective uses weights 20 for the velocity
boundary term, 1 for the scaled gauge term, and 5 for the normalized flux
term. Finite-element interior values are not used in this objective.

\subsection{ForceNet}
The second network has five hidden layers with 90 tanh units and maps
\begin{equation}
 (u,v,p_v,d_p,x_p,y_p,u_p,v_p)\longmapsto(F_x,F_y).
 \label{eq:forcenet}
\end{equation}
The force scale is the displaced-fluid mass times $U_0^2/L_0$. Under this
convention, the dimensionless Stokes force and particle acceleration are
\begin{equation}
 \vect{F}^{*}=\frac{18}{\Rey d_p^2}
 (\vect{u}-\vect{v}),\qquad
 \frac{d\vect{v}}{dt}=\frac{\vect{F}^{*}}{R}.
 \label{eq:force_scaling}
\end{equation}
The neural approximation is trained on the normalized slip response.
Subtracting its evaluation at zero slip from its evaluation at the actual
slip enforces exactly zero force when $\vect{u}=\vect{v}$. The network
therefore approximates a specified drag relation. Its agreement with that
relation is an approximation test, not discovery of a new constitutive law.

\Cref{fig:basic_architecture} gives a visual summary of both networks.
The role of ForceNet is intentionally kept separate from the carrier-field
approximation so that the two errors can be evaluated independently.

\ResultFigure[.96]{02_baseline_architecture}{Baseline neural architecture.
Each row represents a fully connected MLP; the few displayed nodes are
illustrative, while the numbers above hidden layers give their actual
widths. FlowNet supplies velocity and pressure. ForceNet approximates the
normalized Stokes response and uses a zero-slip subtraction. The first
network is constrained by PDE and boundary residuals; the second by an
algebraic drag residual.}{fig:basic_architecture}

\section{Independent finite-element reference}
The reference solves the same incompressible equations using quadratic
velocity and linear pressure Taylor--Hood elements
\citep{taylor1973,gustafsson2020}. A Picard iteration treats the convective
term, and sparse direct solution is used for the linear subproblems.
The finest channel mesh contains 3,026 triangles, 12,530 velocity degrees
of freedom, and 1,620 pressure degrees of freedom.

Three meshes assess sensitivity. \Cref{tab:mesh} shows that the difference
between the two finest velocity fields is \MeshDifference\%. The finest
reference has a nonlinear residual of approximately
$1.60\times10^{-10}$. Its reported divergence norm is not identically zero;
the mixed discretization enforces incompressibility weakly. The mesh
differences are sensitivity indicators, not a formal uncertainty estimate
or a grid-convergence index \citep{roache1998,celik2008}. % Citation-plan C20

\begin{table}[htbp]\centering\small
\caption{Reference-mesh assessment. Velocity differences are measured
relative to the finest field on the comparison samples.}
\label{tab:mesh}
\begin{tabularx}{\textwidth}{@{}rrrrY@{}}\toprule
Triangles & Velocity DOFs & Pressure DOFs & Difference (\%) & Nonlinear residual\\\midrule
\EvidenceRows{mesh}
\bottomrule\end{tabularx}
\end{table}

The first successful training seed is used for the principal comparison;
it is not selected by its error against the reference. The reported final
evaluation restores completed checkpoints. Consequently, the training
budgets associated with those checkpoints, rather than the defaults of a
later evaluation command, identify the computations that produced the
weights. The computational record in \cref{app:reproducibility} preserves
this distinction.

\section{Analytical plane-flow verification}
Before interpreting curved-channel results, I consider plane Poiseuille
flow in a straight channel with half-width $h$. In the baseline units, the
analytical solution is
\begin{equation}
 u=1-(y/h)^2,\qquad v=0,\qquad
 p_v=-\frac{2x}{h^2}+C.
 \label{eq:poiseuille}
\end{equation}
The finite-element relative velocity error is approximately
$8.19\times10^{-10}$. The PINN velocity error is
\PoiseuilleVelocity\%, and its gauge-aligned pressure error is
\PoiseuillePressure\%. The velocity metric includes both components;
\cref{fig:poiseuille} displays the streamwise profile for readability.
Pressure is evaluated on points distributed along the channel, with each
array's mean removed.

The near-exact finite-element result for this polynomial plane flow is
consistent with the approximation spaces and simple geometry. It should
not be interpreted as the error of the curved-channel reference. The
curved problem requires its own mesh-sensitivity assessment.

\ResultFigure{03_poiseuille}{Analytical verification of the baseline flow
solvers. Panel (a) compares the streamwise velocity profile with the
parabolic solution. Panel (b) compares pressure after constant-gauge
alignment on spatially distributed samples. Lines and markers retain the
stored predictions; no fitted correction is applied.}{fig:poiseuille}

\section{Curved-channel carrier-flow results}
\Cref{fig:curved_fields} compares the principal neural field with the
finite-element field. Both solutions recover the curved high-speed core
and the pressure decrease through the channel. Common color limits are
used for corresponding speed and pressure panels. The discrepancy panels
show that agreement is spatially nonuniform; a small global norm should
not be interpreted as uniformly small local error.

The velocity error is \FlowVelocity\%, and the centered pressure error is
\FlowPressure\% on the independent field-comparison samples. The maximum
absolute velocity discrepancy is approximately $0.05401$. These values
support agreement for the specified boundary-value problem. They do not
establish zero equation residual: the separately recorded raw continuity
and momentum RMS values are approximately $0.1781$ and $1.8296$ in the
baseline equation convention. Those residuals, the field errors, and the
flux diagnostic measure different aspects of the approximation.

\ResultFigure{04_curved_fields}{Curved-channel carrier field at $\Rey=1$.
The speed panels use $|\vect{u}|/U_0$ and identical scales; the pressure
panels use centered viscous-scale pressure and identical scales.
Velocity discrepancy is the Euclidean component error, whereas pressure
discrepancy is signed. The fields are evaluated on the same stored mapped
grid; no spanwise solution is implied.}{fig:curved_fields}

\begin{table}[htbp]\centering\small
\caption{Supported baseline numerical comparisons. The relative errors
refer to their stated sample sets and normalizations.}
\label{tab:basic_errors}
\begin{tabularx}{\textwidth}{@{}p{.31\textwidth}Yr@{}}\toprule
Problem & Quantity & Error\\\midrule
\EvidenceRows{basic_errors}
\bottomrule\end{tabularx}
\end{table}

\section{Held-out drag and field comparisons}
The normalized ForceNet error is \DragError\% on the independent
2,048-sample evaluation used for the reported metric.
\Cref{fig:heldout} presents velocity and pressure parity together with a
separate 1,000-sample drag diagnostic. The plotted drag set and the aggregate
2,048-sample test are distinct; they need not have identical numerical
errors. This distinction avoids assigning one test's error to another
test's points.

The drag comparison supports the network's approximation of the retained
creeping-flow law. It cannot establish the accuracy of neglected lift,
history, lubrication, or wall corrections
\citep{maxey1983,goldman1967,brenner1961}. % Citation-plan C21
 Direct analytical Stokes drag
also remains an essential baseline for judging the value of a learned
approximation of this known relation.

\ResultFigure{05_heldout_agreement}{Independent numerical comparisons.
Velocity components and gauge-aligned pressure are compared with the
finite-element reference. The third panel compares normalized learned
drag with the specified Stokes-slip response on the stored 1,000-sample
diagnostic. The dashed lines indicate equality. The aggregate drag error
in \cref{tab:basic_errors} comes from a separate 2,048-sample test.}
{fig:heldout}

\section{Seed variation, reference sensitivity, and flux}
Two independently trained seeds give velocity errors of 1.510\% and
1.429\%, and pressure errors of 0.863\% and 0.886\%, respectively. Their
mean velocity error is \SeedVelocityMean\%, with a sample standard deviation
of \SeedVelocitySD\ percentage points. Two seeds are sufficient to reveal
some training dependence but are insufficient for a strong statistical
uncertainty statement. No confidence interval is inferred.

\begin{table}[htbp]\centering
\caption{Field errors for the two completed training seeds.}
\label{tab:seeds}
\begin{tabular}{@{}rrr@{}}\toprule
Seed & Velocity error (\%) & Pressure error (\%)\\\midrule
\EvidenceRows{seeds}
\bottomrule\end{tabular}
\end{table}

\ResultFigure{06_reference_and_seeds}{Reference and training sensitivity.
The left panel compares the two coarser finite-element meshes with the
finest reference; the right panel shows the two completed training seeds.
The finest mesh is a comparison reference, not an exact solution. Seed
labels identify the individual observations rather than confidence bounds.}
{fig:reference_seeds}

The maximum section-flux deviation is \FluxDeviation\%.
\Cref{fig:flux} shows the full stored section sequence normalized by the
prescribed flux. The variation is finite and should remain visible even
though the study's 5\% flux criterion is met. Together with boundary RMS
of approximately $0.00420$, this result provides an integrated check on
the velocity field that complements the sample-relative error.

\ResultFigure[.95]{07_flux}{Section flux in the curved baseline, normalized
by the prescribed value $Q_0=4h/3=0.2$. The dashed line is the exact
constant-flux target, and every stored PINN section is shown. The maximum
deviation is \FluxDeviation\%; the curve has not been rescaled to enforce
agreement.}{fig:flux}

\section{Particle-transport interpretation}
The baseline evaluation completes all 119 particle or tracer integrations:
49 use PINN with ForceNet, 49 use
FEM with exact Stokes drag, seven use PINN with exact Stokes drag, and
14 are tracer integrations split equally between the PINN and FEM fields.
These counts include the parameter-study cases; the primary endpoint
comparison concerns seven initial conditions. The reference velocity
interpolator differs from direct finite-element evaluation by
\InterpolationDifference\% on its assessment samples. These are useful
computational checks, but neither alone proves accurate particle outcomes.

Among the seven baseline particle initial positions, six endpoint classes
agree with the reference. The largest width-normalized RMS position
discrepancy is approximately 34.97\%, and the lowest common-time coverage
is approximately 73.55\%. Three initial conditions meet all endpoint,
position-error and common-time criteria; six meet the endpoint criterion
alone. Detailed values are retained in
\cref{tab:path_limits} so that the scope of the carrier-flow findings is
unambiguous. The quantitative conclusions of this chapter concern the
carrier field, drag approximation, and the seven tabulated path comparisons.

This outcome is physically and numerically informative. Bulk field
agreement and a good approximation of a specified drag law are necessary
components, but their combination is not enough to certify wall-sensitive
transport. The justified baseline conclusion is that the carrier and drag
components agree well with their stated references, while the full
trajectory prediction remains an improvement target.

The present comparisons do not isolate a single cause for the trajectory
differences. Carrier-field approximation, ForceNet error, reference
interpolation, ODE integration, and event handling are separate candidates.
The recorded integrator uses relative and absolute tolerances
$2\times10^{-6}$ and $2\times10^{-8}$, with RK45 or Radau chosen according
to the response scale. A factorial comparison holding the integrator and
wall-event rule fixed is needed to attribute the discrepancies.

\section{Chapter conclusions}
The baseline establishes a reproducible carrier-flow and drag foundation
for the subsequent coupling study. Analytical plane flow, independent
curved-channel fields, two completed training seeds, and section flux
provide complementary evidence. The results also set a boundary on the
claim: the present field accuracy does not establish reliable deposition
or all baseline particle endpoints. That boundary becomes more important
when reaction forcing and geometric narrowing are introduced.

\chapter{Gaussian Parcels in a Straight Channel}\label{ch:gaussian}

\section{From individual particles to represented populations}
The second stage adds hydrodynamic feedback while retaining a simple
carrier geometry. A parcel represents $N_i$ identical spherical particles
whose centers are approximated by a common planar state. Its mass and
hydrodynamic reaction are weighted by that multiplicity. The Gaussian
distribution describes how the represented force interacts with the
carrier; it does not resolve the surface of a sphere.

The dimensionless fluid domain is $\Omega=[0,3]\times[0,1]$, periodic in
$x$ with stationary no-slip walls at $y=0$ and $y=1$. The channel has no
open inlet or outlet. The initial fluid velocity is
\begin{equation}
 \vect{u}_0(y)=(6y(1-y),0),
 \label{eq:gaussian_initial}
\end{equation}
whose mean streamwise speed is one. A body-force acceleration
$\vect{f}_d=(12/\Rey_H,0)$ maintains this profile in the unladen steady
limit. Parcel initial velocities are zero. The reference depth ratio is
$b=1$, and the evaluated model uses $\Rey_H=20$, $R=1000$, six parcels,
and Gaussian width $\sigma=0.14$.

\section{Coupled equations}
The fluid equations use the inertial pressure convention defined in
\cref{eq:gaussian_scales}:
\begin{align}
 \nabla\cdot\vect{u}&=0,\\
 \partial_t\vect{u}+(\vect{u}\cdot\nabla)\vect{u}
 &=-\nabla p_i+\frac{1}{\Rey_H}\nabla^2\vect{u}
   +\vect{S}+\vect{f}_d.
 \label{eq:gaussian_ns}
\end{align}
The parcel equations are
\begin{equation}
 \dot{\vect{q}}_i=\vect{v}_i,\qquad
 \dot{\vect{v}}_i=\frac{\vect{U}_i-\vect{v}_i}{\St_i}
 +\vect{g}_{\mathrm{eff}},\qquad
 \vect{F}_i=\frac{m_i}{\St_i}(\vect{U}_i-\vect{v}_i).
 \label{eq:gaussian_particle}
\end{equation}
Here $\vect{g}_{\mathrm{eff}}$ is a specified effective acceleration,
including any buoyancy adjustment in the physical mapping. The reaction
source contains the hydrodynamic drag only. Depositing parcel gravity again
as fluid drag would count a different force twice.

\section{Normalized Gaussian transfer}
Let $\vect{x}_a$ denote the fixed transfer-quadrature locations and
$\Delta A_a$ their area weights. A periodized raw Gaussian is evaluated
between $\vect{x}_a$ and each parcel center. In the two-dimensional channel,
its structure is
\begin{equation}
 G_{ai}=\sum_{k\in\mathbb{Z}}
 \exp\left[-\frac{(x_a-q_{ix}+kL_x)^2+(y_a-q_{iy})^2}
 {2\sigma^2}\right].
 \label{eq:raw_kernel}
\end{equation}
The implementation retains a finite image set with distant tails beyond
eight widths after choosing the nearest image. A global Gaussian prefactor
is unnecessary because normalization removes it. Define
\begin{equation}
 Z_i=\sum_aG_{ai}\Delta A_a,\qquad W_{ai}=G_{ai}/Z_i,
 \qquad \sum_aW_{ai}\Delta A_a=1.
 \label{eq:normalized_kernel}
\end{equation}
Both fluid interpolation and force spreading use these weights:
\begin{equation}
 \vect{U}_i=\sum_aW_{ai}\vect{u}_a\Delta A_a,\qquad
 \vect{S}_a=-\sum_iW_{ai}\vect{F}_i.
 \label{eq:shared_transfer}
\end{equation}
Near a wall, $Z_i$ accounts for the actual bounded-domain quadrature.
At an arbitrary collocation point, the numerator uses that point, while
the denominator remains the value from the fixed transfer quadrature.

\subsection{Force and work identities}
Substituting \cref{eq:shared_transfer} into the quadrature sum gives
\begin{equation}
 \sum_a\vect{S}_a\Delta A_a=-\sum_i\vect{F}_i,
 \label{eq:force_identity}
\end{equation}
and
\begin{equation}
 \sum_a\vect{u}_a\cdot\vect{S}_a\Delta A_a
 =-\sum_i\vect{U}_i\cdot\vect{F}_i.
 \label{eq:work_identity}
\end{equation}
The combined interphase contribution to fluid and parcel kinetic-energy
rates is therefore
\begin{equation}
 \mathcal{P}_{\mathrm{drag}}
 =\sum_i(\vect{v}_i-\vect{U}_i)\cdot\vect{F}_i
 =-\sum_i\frac{m_i}{\St_i}|\vect{U}_i-\vect{v}_i|^2\leq0.
 \label{eq:drag_dissipation}
\end{equation}
These are semidiscrete properties of the shared transfer. A trained PINN
with nonzero residual need not satisfy the complete global momentum or
energy balance exactly. External driving can also make total kinetic energy
increase even though the drag contribution is dissipative.

\section{Neural representation and hard constraints}
Two MLPs, each with three hidden layers of 32 tanh units, represent the fluid
correction and parcel histories. The fluid network uses wall-normal position,
time, and three periodic Fourier harmonics in $x$. The ansatz is
\begin{equation}
 \vect{u}(\vect{x},t)=\vect{u}_0(y)
 +\frac{t}{T}y(1-y)\vect{n}_{u}(\vect{x},t).
 \label{eq:fluid_ansatz}
\end{equation}
It enforces the initial velocity and stationary wall velocity. Periodicity
follows from the Fourier features. A time-dependent pressure offset is
removed at $(0,0)$; no independent initial-pressure condition is imposed.

The parcel network uses the relaxation time in its output envelopes:
\begin{align}
 b_i(t)&=1-\exp(-t/\St_i),&
 a_i(t)&=t-\St_i b_i(t),\\
 \vect{q}_i(t)&=\vect{q}_{i0}+t\vect{v}_{i0}
                    +a_i(t)\vect{n}_{q,i}(t),&
 \vect{v}_i(t)&=\vect{v}_{i0}+b_i(t)\vect{n}_{v,i}(t).
 \label{eq:parcel_ansatz}
\end{align}
The envelope imposes the initial states and represents the known response
scale. It supplies no target fluid velocity or target coupled trajectory.
Position and velocity remain separate trainable functions, linked by
$\dot{\vect{q}}_i-\vect{v}_i$ in the objective.

The weighted mean-square residual objective assigns weights 1 to fluid
momentum, 2 to continuity, 5 to parcel kinematics, 1 to parcel acceleration,
and 100 to center confinement. All calculations in this formulation use
float64. \Cref{fig:gaussian_architecture} emphasizes the common quadrature
that couples the representations. The default calculations are
collision-free; independently tested impact primitives are not treated as
a validated collisional suspension model.

\ResultFigure[.96]{08_gaussian_architecture}{Coupled Gaussian formulation.
The two stylized network symbols indicate distinct MLPs, each with three
32-unit hidden layers; their radial arrangement is illustrative and does
not specify literal layer connectivity. Shared quadrature transfers sampled
fluid velocity to parcels and hydrodynamic reaction to the fluid. The
complete objective includes fluid and parcel residuals.}
{fig:gaussian_architecture}

\section{Reduced Stokes-response verification}
The first reduced test, V01, solves the constant-forcing relaxation problem
with dimensionless response time one:
\begin{equation}
 \dot v=1-v,\qquad \dot z=v,\qquad v(0)=z(0)=0.
 \label{eq:stokes_test}
\end{equation}
Its analytical response is
\begin{equation}
 v=1-e^{-t},\qquad z=t-1+e^{-t}.
 \label{eq:stokes_exact}
\end{equation}
The reference arrays are excluded from training and used only after the
physics-constrained fit. The relative velocity and displacement errors are
\VOneVelocity\% and \VOneDistance\%, with independent ODE residual RMS
approximately $1.35\times10^{-3}$. Both the stated relative-error and
residual criteria are satisfied. \Cref{fig:stokes_response} shows the
stored predictions against the analytical response.

\ResultFigure{09_stokes_response}{V01 reduced Stokes response. The PINN
is fitted to the ODE residual and initial conditions; analytical values are
held out. Open markers are subsampled stored predictions for legibility;
the metrics use every evaluation sample. This test verifies the reduced
response equation, not the full two-way channel.}{fig:stokes_response}

\section{Two-way momentum-exchange verification}
The second reduced problem, V02, considers spatially uniform fluid and
particle velocities, fluid mass one, parcel mass loading $\beta=0.25$,
and response time $\St=0.2$. With $u(0)=0$ and $v(0)=1$,
\begin{equation}
 \dot u=-\frac{\beta}{\St}(u-v),\qquad
 \dot v=\frac{u-v}{\St}.
 \label{eq:two_way_test}
\end{equation}
The solution has constant total momentum $u+\beta v=0.25$ and decaying
slip. The relative fluid- and particle-velocity errors are
\VTwoFluid\% and \VTwoParticle\%, respectively. The independent residual
RMS is approximately $6.12\times10^{-3}$.

In this reduced PINN, the output relation $v=1-4u$ imposes total momentum
exactly. Its zero momentum error is consequently a hard-constraint check,
not an independent demonstration that a general coupled solution conserves
momentum. The nontrivial checks are the predicted relaxation and its ODE
residual. \Cref{fig:two_way_response} also shows the kinetic-energy decay,
consistent with the sign in \cref{eq:drag_dissipation}.

\ResultFigure{10_two_way_response}{V02 uniform two-way relaxation.
The two velocities approach a common value while kinetic energy decreases.
The total-momentum line is exact by the imposed output relation. Velocity
agreement and the independently sampled ODE residual assess whether the
learned temporal response satisfies the governing problem.}
{fig:two_way_response}

\section{Kernel and carrier-flow verification}
\subsection{Normalization, overlap, and parcel splitting}
The bounded-domain normalization test V03 probes periodic-seam and wall
locations. Its maximum recorded discrepancy is approximately
$8.44\times10^{-15}$ across the checks included in that test. The displayed
integrals in \cref{fig:kernel_checks} differ from one by only a few units
of machine precision. The larger summary discrepancy also includes the
constant-field sampling check and should not be assigned to an individual
integral in the figure.

V04 compares a three-dimensional product quadrature with an independent
free-space Gaussian overlap array. The relative difference is approximately
$1.29\times10^{-15}$. The CSV column retained from the numerical workflow
uses a generic prediction name; here it is correctly identified as a
quadrature result rather than a learned PINN field. This is a kernel-integral
test and does not establish a three-dimensional channel solution.

V09 splits a fixed population of 10,000 particles into 10, 50, and 100
coincident parcels. The maximum force-density difference is approximately
$2.66\times10^{-15}$. Because the split centers coincide and total mass
is unchanged, this tests deterministic representation invariance. It is not
a convergence study of independently dispersed parcel populations.

\ResultFigure{12_kernel_verification}{Gaussian transfer checks. Panel (a)
shows V04's independent three-dimensional overlap integral and product
quadrature. Panel (b) shows the stored V03 normalization values divided by
machine precision. The kernels are verified on their declared quadrature;
neither panel is a flow-field prediction.}{fig:kernel_checks}

\subsection{Unsteady incompressible flow}
V10 considers the periodic two-dimensional Taylor--Green vortex with
viscosity $\nu=0.01$ over $0\leq t\leq5$. The exact velocity is
\begin{equation}
 u=\sin x\cos y\,e^{-2\nu t},\qquad
 v=-\cos x\sin y\,e^{-2\nu t},
 \label{eq:tg_velocity}
\end{equation}
and the gauge-aligned pressure is
\begin{equation}
 p=\frac{1}{4}(\cos2x+\cos2y-2)e^{-4\nu t}.
 \label{eq:tg_pressure}
\end{equation}
The initial field is imposed, while the analytical temporal decay is held
out of training. Independent sample-relative errors are \TGU\% for $u$,
\TGV\% for $v$, and \TGP\% for pressure. The independent residual RMS is
approximately $7.38\times10^{-4}$. These results support the unsteady
carrier implementation for this smooth periodic test.

\ResultFigure{11_taylor_green}{V10 Taylor--Green verification on independent
space--time samples. Each panel compares the predicted quantity with its
analytical value, using the same pressure gauge. The test contains no
parcels and establishes carrier-flow verification rather than coupled
particle-transport validation.}{fig:taylor_green}

\begin{table}[htbp]\centering\small
\caption{Accepted solution-observable errors in the reduced Gaussian-model
verification problems. Residual criteria are checked separately.}
\label{tab:gaussian_errors}
\begin{tabularx}{\textwidth}{@{}lYr@{}}\toprule
Test & Observable & Relative error\\\midrule
\EvidenceRows{gaussian_errors}
\bottomrule\end{tabularx}
\end{table}

\section{Status of the coupled straight-channel calculations}
The evaluated coupled channel uses one training seed and four combinations
of Stokes number and mass loading. The training schedule is 1,200 Adam
steps followed by a 200-iteration L-BFGS budget; the transfer quadrature
is $16\times16$. All four calculations produce finite outputs, but none
meets the complete residual screen. \Cref{tab:coupled_scope} reports the
momentum and parcel-acceleration RMS values that bound their interpretation.

\begin{table}[htbp]\centering\small
\caption{Coupled straight-channel screening. Every required held-out
residual must be below 0.03; the momentum residual exceeds that bound in
all four cases. These are diagnostic calculations, not accepted transport
predictions.}\label{tab:coupled_scope}
\begin{tabularx}{\textwidth}{@{}rrrrY@{}}\toprule
$\St$ & $\beta$ & Momentum RMS & Parcel RMS & Complete screen\\\midrule
\EvidenceRows{coupled_scope}
\bottomrule\end{tabularx}
\end{table}

The lowest listed momentum RMS is $0.0407$, still above the declared bound.
The full screen also requires all centers to remain inside the channel,
a trained finite checkpoint without training error, local volume fraction
below 0.01, and slip Reynolds number no greater than one. The observed
maximum slip Reynolds numbers range from approximately $0.234$ to $0.642$.
These satisfy that loose screening bound but do not establish the
asymptotically small-slip regime required for accurate Stokes drag.
Passing the reduced tests therefore does not justify selecting a visually
smooth coupled-channel field as a validated result. I exclude these
predictions from quantitative claims about loading effects, sedimentation,
or deposition. The physical model also samples the disturbed fluid without
an undisturbed-velocity correction, so successful future residual screening
would still require independent assessment of the drag closure.

All ten component tests in the evaluated record pass their individual
criteria. Tests V05--V08 concern impact algebra and a frozen-rate sampling
primitive, which are outside the collision-free channel study and are not
used to validate its dynamics. The literature catalogue likewise supplies
no matched physical-validation result for the Gaussian channel. These
scope limits preserve the distinction between tested components and an
accepted coupled prediction.

\section{Chapter conclusions}
The Gaussian formulation establishes exact transfer-quadrature identities
and successful reduced tests for relaxation, kernel normalization, and
unsteady carrier flow. The strongest supported results are these
independently evaluated components. Further optimization, transfer-grid
and time-sampling refinement, multiple seeds, and a matched independent
coupled reference are required before interpreting the fully coupled
channel quantitatively. The next chapter therefore formulates a controlled
constricted-channel extension rather than extending unsupported transport
claims.

\chapter{Proposed Constricted-Channel Study}\label{ch:constriction}

\section{Purpose and status of the extension}
A constriction introduces a varying gap, a changing cross-sectional mean
velocity at fixed flux, and stronger sensitivity to wall treatment.
It provides a natural next problem after the constant-width baseline and
the Gaussian coupling study. It also introduces new requirements that
cannot be met by merely changing the appearance of a straight-channel
figure.

This chapter defines the geometry, governing problem, and evaluation design
for that extension. The geometry drawings are mathematical illustrations.
No constricted-channel velocity, pressure, particle trajectory, or deposition
rate has been computed for the accepted result set. The proposed parameter
values below are study-design choices rather than measured or completed
results.

The obstruction is static and represented by the walls. A progressively
growing particle deposit, moving obstacle, or porous plug would require
additional evolution laws and is outside this design. I use the term
\emph{constriction} for the prescribed narrowing to keep this geometric
effect separate from dynamic clogging.

\section{Parameterized geometry}
Let the nominal channel height be $H$ and let the walls be
$y=\pm h(x)$. A compact, smooth narrowing centered at $x_c$ is defined by
\begin{equation}
 h(x)=\frac{H}{2}[1-\delta f(\xi)],\qquad
 \xi=\frac{x-x_c}{\ell},\qquad 0\leq\delta<1,
 \label{eq:constriction_height}
\end{equation}
where
\begin{equation}
 f(\xi)=
 \begin{cases}
 \left[\dfrac{1+\cos(\pi\xi)}{2}\right]^2,&|\xi|<1,\\[2mm]
 0,&|\xi|\geq1.
 \end{cases}
 \label{eq:constriction_profile}
\end{equation}
The profile and its first three derivatives match the constant-height
segments at the joins. Its minimum gap is $H_t=(1-\delta)H$. The initial
design uses $L_x/H=8$, $x_c/H=4$, and $\ell/H=1.2$, with nominal severities
$\delta\in\{0,0.2,0.4\}$. These values define a controlled sequence from
the unblocked channel to moderate narrowing. The actual reference study
must check that upstream and downstream lengths are sufficient.

\Cref{fig:constriction_geometry} displays the geometry and height profiles.
Only geometric quantities are plotted. In particular, the height curves
must not be interpreted as velocity or pressure results.

\ResultFigure{13_constriction_geometry}{Proposed constriction. Panel (a)
defines the throat, characteristic length, and reference sections. Panel
(b) shows the mathematical height profiles for the proposed severity
values. The illustrated channel has $L_x/H=8$ and $\ell/H=1.2$; no fluid
or particle solution is represented.}{fig:constriction_geometry}

\section{Carrier-flow boundary conditions}
For an open-channel extension, the governing equations retain the inertial
pressure scale and the same Newtonian constitutive law. I prescribe a
parabolic inlet velocity with a specified flux $Q$ and no slip on both
walls. A downstream traction condition provides an outlet reference rather
than imposing a second velocity profile. The finite-element reference and
the PINN must implement exactly the same outlet condition and pressure
convention.

This is a different boundary-value problem from both earlier stages.
The curved baseline has prescribed velocities at both ends, whereas the
straight Gaussian channel is periodic and body-force driven. Comparisons
across these stages must therefore identify which changes follow from
geometry and which follow from the drive or boundary conditions. A useful
control is a straight open channel with the same inlet and outlet treatment
as the constricted case.

For unsteady particle introduction, the initial carrier field should be an
independently converged unladen solution in the constricted domain.
Imposing the original straight-channel profile throughout the narrowed
region would violate the new wall geometry. An alternative start from rest
requires a compatible smooth inlet ramp and its own transient verification.
The primary design uses the steady unladen field to avoid mixing inlet
startup with particle response.

\section{Geometry-aware representation}
The fixed-width centerline mapping of the baseline cannot represent this
variable-height domain without modification. A convenient mapping is
\begin{equation}
 (x,\eta)\longmapsto(x,y)=(x,h(x)\eta),\qquad -1\leq\eta\leq1.
 \label{eq:constriction_map}
\end{equation}
Its area Jacobian is $h(x)$. If the differential residual is evaluated in
mapped coordinates, the chain rule must include this geometry dependence.
For example,
\begin{equation}
 \left.\partial_x\right|_y
 =\left.\partial_x\right|_\eta
 -\eta\frac{h'(x)}{h(x)}\partial_\eta,\qquad
 \partial_y=\frac{1}{h(x)}\partial_\eta.
 \label{eq:mapped_derivatives}
\end{equation}
Evaluating automatic derivatives directly with respect to physical
coordinates is another option. Either route requires a manufactured
derivative test before training a new flow problem.

Periodic Fourier features in the streamwise direction would impose the
wrong structure for the proposed open domain. They must be replaced by
nonperiodic coordinate features. A wall-distance factor can enforce the
stationary wall velocity, but it must be consistent with the nonzero inlet
condition and with the mapped derivatives \citep{sukumar2022}. % Citation-plan C27
 A divergence-free streamfunction
ansatz is a possible alternative, provided that its boundary and pressure
conditions are verified. These are implementation choices to evaluate,
rather than capabilities already established by the straight-channel model.

\section{Gaussian support near the narrowing}
The force-transfer rule must use quadrature over the fluid domain:
\begin{equation}
 Z_i=\sum_{a\in\Omega_\delta}
 G(\vect{x}_a-\vect{q}_i)\Delta A_a,\qquad
 W_{ai}=\frac{G(\vect{x}_a-\vect{q}_i)}{Z_i}.
 \label{eq:constriction_normalization}
\end{equation}
For mapped quadrature, $\Delta A_a$ includes the Jacobian from
\cref{eq:constriction_map}. The same points, weights, and normalizer must
be used for interpolation and reaction-force deposition. Under that
condition, the derivations of \cref{eq:force_identity,eq:work_identity}
remain valid on the new quadrature.

\Cref{fig:boundary_kernel} illustrates why normalization must follow the
actual domain. Some Gaussian support intersects the wall region near the
throat. Removing those points without updating $Z_i$ would lose force.
Renormalizing restores the integrated force, but changes the shape of the
sampling operator, an issue related to coarse-graining choices
\citep{sun2015}. Physical wall corrections \citep{goldman1967,brenner1961}
and undisturbed-velocity recovery \citep{horwitz2016,ireland2017} remain
separate modeling requirements. % Citation-plan C28

\ResultFigure[.95]{14_boundary_kernel}{Illustration of a Gaussian centered
near the upper throat wall. Shading is a prescribed kernel shape restricted
to the proposed fluid domain, not a computed concentration or force field.
The normalizer must use the same fluid-domain quadrature as both transfer
operations.}{fig:boundary_kernel}

Three ratios require explicit monitoring:
\begin{equation}
 \frac{d}{H_t},\qquad \frac{\sigma}{H_t},\qquad
 \frac{\sigma}{\Delta_{\mathrm{quad}}},
 \label{eq:constriction_ratios}
\end{equation}
where $\Delta_{\mathrm{quad}}$ denotes a local quadrature-spacing measure.
The first concerns physical confinement; the second concerns regularization
relative to the gap; the third concerns the numerical resolution of that
regularization. They address different questions and should be varied
separately, with coarse-graining and transfer-resolution studies providing
related context \citep{sun2015,capecelatro2013}. % Citation-plan C29
 If physical wall effects dominate, the present unbounded
Stokes closure may need replacement before more accurate training is useful.

\section{Particle and parcel wall conditions}
An individual sphere has an admissible center region obtained by offsetting
the wall by its radius. For a sloped wall, the offset is based on signed
normal distance to the boundary, not only the vertical gap. A contact event
must therefore use the actual geometry. The baseline's irreversible
sticking rule can be retained as a controlled modeling choice, but it must
not be called a prediction of deposition physics.

A parcel introduces an additional issue: its center represents many
particles. Stopping the entire represented population at the first center
contact may create a strong and artificial capture rule. The proposed
first study therefore focuses on carrier fields and freely transported
dilute parcels whose centers remain inside the admissible region. A
deposition extension should specify whether it resolves individual
particles, splits parcels, or uses a separately justified population-level
wall model. No capture efficiency is assigned before making that choice.

\section{Controlled evaluation sequence}
The extension uses staged verification consistent with the distinction
between solving equations and establishing physical validity
\citep{oberkampf2002}. The following sequence is the study design adopted
here. % Citation-plan C30
\begin{enumerate}
\item Recover the straight, unladen open-channel solution at $\delta=0$
with the new inlet and outlet conditions. Verify pressure gauge, wall
velocity, and section flux before adding particles.
\item Solve the unladen constricted channel with an independent mesh-based
method. Refine the throat and the upstream and downstream transition
regions; monitor pressure drop and the complete velocity field.
\item Introduce individual one-way particles with identical initial
conditions in the neural and reference fields. Compare position,
common-time coverage, and endpoint class separately.
\item Add a single Gaussian parcel at low loading and verify force/work
transfer on the irregular-domain quadrature. Compare the coupled fluid
response with an independently discretized regularized model.
\item Vary severity and loading only after the earlier controls pass.
Keep physical diameter, total represented mass, and kernel width explicit.
\end{enumerate}

The initial proposed Reynolds numbers are $1$, $5$, and $20$ at nominal
height. These are design points rather than an observed regime map.
Particle response times should be selected to retain an acceptably small
slip Reynolds number at the throat. Mass-loading changes must be made
through represented population while holding diameter fixed when the
purpose is to isolate feedback. A separate diameter sweep changes both
response time and confinement, and should be interpreted accordingly.

\section{Observables and acceptance criteria}
The primary observables are the section flux, throat velocity profile,
pressure drop between fixed reference sections, and field error against
the matched reference. For particles, residence time and endpoint class
are added only after trajectory agreement is demonstrated. The reference
sections must remain identical across severity values so that the pressure
drop represents the same geometric interval.

At fixed flux, incompressibility gives the exact sectional mean relation
\begin{equation}
 \overline{u}_x(x)=\frac{Q}{2h(x)},\qquad
 \frac{\overline{u}_{x,t}}{\overline{u}_{x,0}}
 =\frac{1}{1-\delta}.
 \label{eq:mean_acceleration}
\end{equation}
This is a geometric conservation implication, not a computed centerline
velocity or a prediction of recirculation. A lubrication estimate,
derived in \cref{app:derivations}, gives a useful additional check only
when its small-slope and low-inertia assumptions apply.

\begin{table}[htbp]\centering\small
\caption{Proposed acceptance protocol for the constricted study. Limits
must be fixed before interpreting the associated parameter sweep; these
checks have not been executed.}\label{tab:constriction_protocol}
\begin{tabularx}{\textwidth}{@{}p{.26\textwidth}YY@{}}\toprule
Quantity & Proposed check & Interpretation\\\midrule
Unladen field & $E_u,E_p<5\%$ against a matched reference & Carrier-field agreement\\
Reference refinement & Changes in $\Delta p$ and field norms below 2\% & Numerical sensitivity, not exact uncertainty\\
Section flux & Maximum relative variation below 2\% & Integrated mass balance\\
Transfer & Force/work discrepancies consistent with roundoff & Algebraic operator consistency\\
Particle paths & RMS position below $0.05W$; common interval above 95\%; all endpoint classes matched & Transport accuracy for stated initial conditions\\
Physical regime & Report maximum $\Rey_p$, $d/H_t$, local $\alpha$ and wall distance & Closure applicability; no universal threshold inferred\\
\bottomrule\end{tabularx}
\end{table}

The error limits in \cref{tab:constriction_protocol} are proposed numerical
screens. In addition, residuals should be reported by component on a
separate set of throat, near-wall, and bulk points. A global average must
not conceal a concentration of error near the narrowing. The Gaussian
width study should remain separate from spatial and temporal refinement,
and multiple seeds should be included before drawing a quantitative trend.

\section{Expected mechanisms and testable statements}
Conservation at fixed flux predicts a larger section-mean speed as the
gap narrows. The accompanying changes in particle transit depend on
response time, local flow, and the retained force law. I treat changes in
residence time, wall approach, and downstream recovery as hypotheses for
the controlled study. Recirculation, enhanced deposition, or a particular
pressure-loss factor cannot be inferred from the geometry alone.

At fixed imposed pressure difference, flux becomes an outcome instead.
That experiment addresses a different question and must not be pooled with
the fixed-flux design. Keeping these operating conditions separate is
essential for a meaningful comparison of unblocked and constricted cases.

\section{Chapter conclusions}
The constricted-channel stage is defined at the study-design level through
its geometry, boundary-condition choices, transfer consistency, and
evaluation requirements. The outlet traction and its reference value must
be fixed before executing the matched solvers.
It extends the study in a controlled direction while preserving the
distinction between mathematical consequences, proposed experiments, and
completed numerical findings. Its next milestone is a verified unladen
reference problem, followed by one-way transport and then two-way Gaussian
feedback.

\chapter{Discussion}\label{ch:discussion}

\section{What the baseline establishes}
The baseline demonstrates that a physics-informed field representation can
agree closely with an independently discretized curved-channel solution
for the specified steady, low-Reynolds-number problem. The velocity and
pressure errors are below 2\% and 1\%, respectively, while the drag-network
error is below 1\%. The analytical plane-flow test provides an additional
check that is independent of the curved-channel finite-element field.
Reference-mesh sensitivity and section flux add evidence beyond a single
parity plot.

The trajectory comparison evaluates an additional quantity beyond field
agreement. Six of seven endpoint classes agree, while three of seven
trajectories meet the complete position, common-time and endpoint criteria.
A particle integrates the local field over time; force approximation,
interpolation and event handling also affect its result. Their individual
contributions require a controlled comparison. These findings motivate
direct transport assessment alongside the carrier-field metrics.

The force network also raises a methodological question. Since the
retained Stokes law is known analytically, its neural approximation is
most useful as a controlled component of a broader learning framework.
The present error does not show that the learned force is superior to the
direct analytical formula. A future assessment should compare the two
within identical carrier fields and identify whether the approximation
adds cost, flexibility, or error.

\section{What conservative transfer establishes}
The shared Gaussian operator gives a stronger result of a different kind.
The force and work identities follow from the normalization and the common
quadrature. Their validity does not depend on a low training loss. This
property removes one avoidable source of inconsistent interphase exchange
and remains available when the carrier geometry changes, provided that the
same fluid-domain quadrature is used consistently.

The identities do not guarantee the accuracy of the coupled solution.
The carrier momentum equation can retain a residual even when the sum of
deposited forces is exact. In the evaluated coupled channel, the remaining
momentum residual exceeds the study's bound in every tested parameter
combination. The correct inference is that the operator design is verified
while the quantitative coupled-channel prediction remains unresolved.

The successful reduced tests address important parts of the calculation:
response-time dynamics, reciprocal momentum exchange, kernel integrals,
and unsteady incompressible flow. Each isolates a necessary capability.
Their combination is a verification chain, not a substitute for a matched
independent two-way channel reference. This distinction is consistent with
the separation between numerical verification and physical validation
\citep{roache1998}.

\section{Optimization, residuals, and independent error}
The completed calculations illustrate why optimization metrics and
independent solution errors should be reported together. Residual
normalization changes their numerical magnitude, and a weighted objective
can favor some components over others. Reported raw residuals should
therefore retain their equation convention. In particular, comparing the
baseline momentum RMS directly with the Gaussian momentum RMS would
ignore different scaling, geometry, and transient content.

Training difficulties in physics-informed learning are well documented
\citep{wang2021,wang2022,krishnapriyan2021}. They motivate focused improvements such as % Citation-plan C31
component scaling and more informative collocation. Residual-based
adaptive sampling \citep{wu2023} and adaptive loss weighting
\citep{mcclenny2023} are options for a future controlled comparison.
A further option is better temporal % Citation-plan C32
resolution of the response scale, including causality-aware training
for the transient coupled problem \citep{wang2024}. These approaches were
not part of the recorded runs. Their value must be assessed % Citation-plan C33
using the same independent checks and the original acceptance criteria.

\section{Regularization and model validity}
The Gaussian width controls how the represented force is distributed.
It affects both the sampled fluid velocity and the carrier response.
A smaller width is not automatically more physically accurate, especially
if the transfer quadrature underresolves it. A larger width can smooth
localized forcing while changing the regularized interaction. Width
sensitivity must therefore be distinguished from convergence at fixed
model parameters.

The undisturbed-velocity issue is equally important. The current parcels
sample the disturbed fluid directly \citep{horwitz2016,ireland2017}. % Citation-plan C34
 A reduced Stokes-response test cannot
resolve the resulting self-interaction question in a spatially coupled
flow. Comparison with an independent solver of the same regularized
equations would test numerical consistency; assessment against a suitable
physical reference would additionally test the closure.

Confinement further separates numerical and physical accuracy. A model
can converge accurately to equations that omit important wall effects.
For that reason, the proposed constriction study records particle-to-gap
ratio and local slip Reynolds number in addition to ordinary solver errors.
If those quantities indicate a breakdown of the retained approximation,
more training is not the appropriate physical correction.

\section{Transferability across the three stages}
\Cref{tab:claim_scope} summarizes the actual scope of the findings. The
baseline establishes a carrier-field comparison for one geometry and
Reynolds number, with two seeds. The Gaussian tests establish specific
reduced dynamics and transfer properties. The constriction stage provides
a defined study design. None of these establishes a universal neural
operator over channel shape, loading, and Reynolds number.

\begin{table}[htbp]\centering\small
\caption{Supported conclusions and their boundaries.}\label{tab:claim_scope}
\begin{tabularx}{\textwidth}{@{}p{.24\textwidth}YY@{}}\toprule
Stage & Supported statement & Remaining boundary\\\midrule
Baseline carrier & Analytical and independent numerical field agreement & One declared geometry/Re case; finite residuals\\
Baseline particles & Completed integrations and drag approximation & Six of seven matching endpoint classes; three complete path comparisons within criteria\\
Gaussian transfer & Discrete force/work consistency and component verification & Disturbed-velocity closure and coupled field error\\
Straight coupled channel & Executed diagnostic predictions & Momentum RMS $0.0407$--$0.3854$ against a $0.03$ criterion\\
Constriction & Defined geometry and evaluation protocol & No completed constricted-channel solution\\
\bottomrule\end{tabularx}
\end{table}

The transferable parts are the nondimensional bookkeeping, reference
matching, shared transfer, and evaluation logic. New physical and geometric
claims require their own executed evidence. In particular, a parameter
slot in a network does not establish learned parametric generalization,
and a geometric illustration does not establish a flow solution.

\section{Implications for the next computational study}
The most informative next work is narrower than a large parameter sweep.
First, the baseline trajectory limitation should be isolated using the
same carrier field with exact and learned drag. Second, the Gaussian
channel should be compared with an independent regularized solver at fixed
width and well-resolved quadrature. Third, the unladen constricted problem
should be solved before introducing moving populations.

This ordering separates carrier approximation, irregular-domain transfer,
wall geometry, particle closure, and optimization. Establishing the controls
in sequence allows each additional mechanism to be interpreted against a
known numerical foundation.

\section{Relationship to earlier methods}
The present contribution is a continuous-time neural representation with
shared, normalized force transfer and explicitly separated field and
trajectory assessment. The transfer structure is related to
Eulerian--Lagrangian methods \citep{peskin2002,capecelatro2013}, while the
fluid representation belongs to the velocity--pressure PINN framework
\citep{jin2021}. Undisturbed-velocity recovery is a distinct modeling
development addressed by Horwitz and Mani \citep{horwitz2016} and Ireland
and Desjardins \citep{ireland2017}. These
relationships identify the design choices adopted here; the reported
comparisons establish performance only for the stated tests and do not
demonstrate general superiority over those earlier methods.

\chapter{Conclusions and Future Work}\label{ch:conclusions}

\section{Principal conclusions}
This thesis develops a staged framework for physics-informed
fluid--particle transport. The progression begins with a one-way coupled
particle problem, introduces Gaussian parcel reaction in a straight
channel, and defines a smooth constriction for the next study. The
completed evidence supports the following conclusions.

First, the baseline carrier network agrees with an independent
finite-element solution of the stated curved-channel problem, with velocity
and gauge-aligned pressure errors of \FlowVelocity\% and
\FlowPressure\%. Analytical plane-flow verification and reference-mesh
sensitivity support this comparison. The normalized learned Stokes
response has an independent error of \DragError\%.

Second, carrier-field agreement does not by itself establish accurate
particle transport. The seven baseline comparisons do not all meet the
endpoint, position, and time-coverage conditions. Deposition and
wall-sensitive parameter trends are therefore outside the accepted
quantitative conclusions.

Third, a normalized Gaussian used for both interpolation and spreading
preserves discrete force and work on the transfer quadrature. The resulting
interphase drag power is nonpositive. These identities are algebraic
properties of the discretization and do not imply that an imperfectly
trained field satisfies every global balance.

Fourth, the reduced Gaussian-model tests recover Stokes response, uniform
two-way relaxation, kernel integrals, and Taylor--Green decay with the
reported errors and residuals. The accepted solution-observable errors
in the listed relaxation and vortex tests are below 0.25\%. The fully
coupled straight-channel predictions remain diagnostic because their
complete residual criteria are not met.

Finally, the constricted-channel extension has a consistent geometric and
evaluation definition, including physical-coordinate derivatives,
fluid-domain kernel normalization, and independent throat-resolution
checks. It has no completed predictive result in the present study.

\section{Contributions}
The study contributes a common evaluation framework connecting carrier
fields, drag approximation, transfer consistency, and particle outcomes.
It gives a dimensionally consistent planar sphere-parcel formulation and
explicitly separates a Gaussian force distribution from a resolved solid
body. It also defines which verified components can be retained when
extending the model to a variable-height channel.

These contributions provide a practical basis for further computation
without assigning unsupported accuracy to complex coupled predictions.
The numerical record and portable manuscript sources preserve the arrays,
criteria, and figure provenance needed to reassess the stated findings.

\section{Priorities for future work}
The immediate priority is to improve and independently verify the
transport-sensitive parts of the calculation. For the baseline, this
means near-wall field assessment and a direct exact-drag comparison. For
Gaussian coupling, it means a matched independent regularized channel
solution, transfer-quadrature refinement, temporal sampling refinement,
and multiple training seeds at fixed physical parameters. Sampling and
numerical-uncertainty methodologies provide a basis for planning those
comparisons \citep{wu2023,celik2008,oberkampf2002}. % Citation-plan C35

The constricted study should then begin with an unladen flow and a
straight-channel boundary-condition control. One-way particles can be
introduced after the carrier reference passes, followed by low-loading
Gaussian feedback. Physical deposition or dynamic clogging requires a
separately specified wall and population model. Experimental comparison
will require matched geometry, operating conditions, particle properties,
and observables before it can establish physical validity.

\clearpage
\renewcommand{\bibname}{References}
\phantomsection\addcontentsline{toc}{chapter}{References}
\begin{thebibliography}{99}
\bibitem{raissi2019}
M. Raissi, P. Perdikaris, and G. E. Karniadakis,
``Physics-informed neural networks: A deep learning framework for solving
forward and inverse problems involving nonlinear partial differential
equations,'' \emph{Journal of Computational Physics}, vol. 378,
pp. 686--707, 2019.
\href{https://doi.org/10.1016/j.jcp.2018.10.045}{doi:10.1016/j.jcp.2018.10.045}.

\bibitem{karniadakis2021}
G. E. Karniadakis, I. G. Kevrekidis, L. Lu, P. Perdikaris, S. Wang, and
L. Yang, ``Physics-informed machine learning,'' \emph{Nature Reviews
Physics}, vol. 3, pp. 422--440, 2021.
\href{https://doi.org/10.1038/s42254-021-00314-5}{doi:10.1038/s42254-021-00314-5}.

\bibitem{wang2021}
S. Wang, Y. Teng, and P. Perdikaris, ``Understanding and mitigating gradient
flow pathologies in physics-informed neural networks,'' \emph{SIAM Journal
on Scientific Computing}, vol. 43, no. 5, pp. A3055--A3081, 2021.
\href{https://doi.org/10.1137/20M1318043}{doi:10.1137/20M1318043}.

\bibitem{wang2022}
S. Wang, X. Yu, and P. Perdikaris, ``When and why PINNs fail to train:
A neural tangent kernel perspective,'' \emph{Journal of Computational
Physics}, vol. 449, article 110768, 2022.
\href{https://doi.org/10.1016/j.jcp.2021.110768}{doi:10.1016/j.jcp.2021.110768}.

\bibitem{maxey1983}
M. R. Maxey and J. J. Riley, ``Equation of motion for a small rigid sphere
in a nonuniform flow,'' \emph{The Physics of Fluids}, vol. 26,
pp. 883--889, 1983.
\href{https://doi.org/10.1063/1.864230}{doi:10.1063/1.864230}.

\bibitem{clift1978}
R. Clift, J. R. Grace, and M. E. Weber, \emph{Bubbles, Drops, and
Particles}. New York: Academic Press, 1978.

\bibitem{gualtieri2015}
P. Gualtieri, F. Picano, G. Sardina, and C. M. Casciola,
``Exact regularized point particle method for multiphase flows in the
two-way coupling regime,'' \emph{Journal of Fluid Mechanics}, vol. 773,
pp. 520--561, 2015.
\href{https://doi.org/10.1017/jfm.2015.258}{doi:10.1017/jfm.2015.258}.
Preprint: \url{https://arxiv.org/abs/1405.6969}.

\bibitem{uhlmann2005}
M. Uhlmann, ``An immersed boundary method with direct forcing for the
simulation of particulate flows,'' \emph{Journal of Computational Physics},
vol. 209, no. 2, pp. 448--476, 2005.
\href{https://doi.org/10.1016/j.jcp.2005.03.017}{doi:10.1016/j.jcp.2005.03.017}.

\bibitem{taylor1973}
C. Taylor and P. Hood, ``A numerical solution of the Navier--Stokes equations
using the finite element technique,'' \emph{Computers \& Fluids}, vol. 1,
no. 1, pp. 73--100, 1973.
\href{https://doi.org/10.1016/0045-7930(73)90027-3}{doi:10.1016/0045-7930(73)90027-3}.

\bibitem{gustafsson2020}
T. Gustafsson and G. D. McBain, ``scikit-fem: A Python package for finite
element assembly,'' \emph{Journal of Open Source Software}, vol. 5,
no. 52, article 2369, 2020.
\href{https://doi.org/10.21105/joss.02369}{doi:10.21105/joss.02369}.

\bibitem{roache1998}
P. J. Roache, \emph{Verification and Validation in Computational Science
and Engineering}. Albuquerque, New Mexico: Hermosa Publishers, 1998.

\bibitem{ferziger2002}
J. H. Ferziger and M. Peri\'c, \emph{Computational Methods for Fluid
Dynamics}, 3rd ed. Berlin: Springer, 2002.
\href{https://doi.org/10.1007/978-3-642-56026-2}{doi:10.1007/978-3-642-56026-2}.

\bibitem{batchelor1967}
G. K. Batchelor, \emph{An Introduction to Fluid Dynamics}.
Cambridge: Cambridge University Press, 1967.

\bibitem{kingma2015}
D. P. Kingma and J. Ba, ``Adam: A method for stochastic optimization,''
in \emph{Proceedings of the 3rd International Conference on Learning
Representations}, 2015.
\url{https://arxiv.org/abs/1412.6980}.

\bibitem{liu1989}
D. C. Liu and J. Nocedal, ``On the limited memory BFGS method for large
scale optimization,'' \emph{Mathematical Programming}, vol. 45,
pp. 503--528, 1989.
\href{https://doi.org/10.1007/BF01589116}{doi:10.1007/BF01589116}.
\end{thebibliography}

\appendix
\chapter{Derivations and Transport Limits}\label{app:derivations}

\section{Uniform two-way relaxation}
For the uniform system in \cref{eq:two_way_test}, define slip $w=v-u$.
Adding the fluid and population-weighted particle equations gives
\begin{equation}
 \frac{d}{dt}(u+\beta v)=0.
\end{equation}
Subtracting the equations gives
\begin{equation}
 \dot w=-\frac{1+\beta}{\St}w,\qquad
 w(t)=w_0\exp\left[-\frac{1+\beta}{\St}t\right].
\end{equation}
If $M_0=u_0+\beta v_0$ is the conserved momentum, the velocities are
\begin{equation}
 u(t)=\frac{M_0-\beta w(t)}{1+\beta},\qquad
 v(t)=\frac{M_0+w(t)}{1+\beta}.
\end{equation}
For $u_0=0$, $v_0=1$ and $\beta=0.25$, both velocities tend to $0.2$.
The total kinetic energy is
\begin{equation}
 K=\frac{1}{2}u^2+\frac{\beta}{2}v^2,
 \qquad \dot K=-\frac{\beta}{\St}(v-u)^2.
\end{equation}
This gives the dissipative reference behavior of V02. It concerns a
closed, uniform exchange system; the driven channel has additional work
and boundary terms.

\section{Adjoint transfer and kinetic energy}
Let $J$ denote the interpolation operator with entries
$J_{ia}=W_{ai}\Delta A_a$. With quadrature mass matrix
$M_A=\operatorname{diag}(\Delta A_a)$, the spreading operator is
$-M_A^{-1}J^{\mathsf T}$ when the parcel force vector stores integrated
forces. Thus
\begin{equation}
 \vect{U}=J\vect{u},\qquad
 \vect{S}=-M_A^{-1}J^{\mathsf T}\vect{F}.
\end{equation}
The fluid work is
$\vect{u}^{\mathsf T}M_A\vect{S}
=-(J\vect{u})^{\mathsf T}\vect{F}$.
Adding parcel drag work produces
\cref{eq:drag_dissipation}. This matrix statement makes explicit why
interpolation and spreading must use the same quadrature and weights.

For a collision-free driven channel with no-slip stationary walls and
periodic streamwise faces, a formally exact solution obeys
\begin{equation}
 \begin{aligned}
 \frac{d}{dt}\left[\frac{1}{2}\int_\Omega|\vect{u}|^2\,dA
 +\sum_i\frac{m_i}{2}|\vect{v}_i|^2\right]
 &=\int_\Omega\vect{u}\cdot\vect{f}_d\,dA
 +\sum_i m_i\vect{v}_i\cdot\vect{g}_{\mathrm{eff}}\\
 &\quad-\frac{1}{\Rey_H}\int_\Omega|\nabla\vect{u}|^2\,dA
 -\sum_i\frac{m_i}{\St_i}|\vect{U}_i-\vect{v}_i|^2.
 \end{aligned}
 \label{eq:energy_budget}
\end{equation}
Pressure and periodic-face contributions cancel under the stated
conditions. The wall velocity is zero, so stationary walls do no work.
Quadrature and neural equation residuals can create a nonzero discrepancy
when this identity is evaluated numerically. Wall tractions still enter
the momentum budget; they must not be omitted merely because wall work
vanishes in the energy budget.

\section{A limiting pressure-drop estimate}
For slowly varying plane walls and negligible inertia, the local
lubrication approximation gives the dimensional velocity
\begin{equation}
 u(x,y)=-\frac{1}{2\mu}\frac{dp}{dx}[h(x)^2-y^2].
\end{equation}
Integrating over $-h\leq y\leq h$ yields
\begin{equation}
 Q=-\frac{2h(x)^3}{3\mu}\frac{dp}{dx},\qquad
 \Delta p=\frac{3\mu Q}{2}\int_{x_a}^{x_b}\frac{dx}{h(x)^3}.
 \label{eq:lubrication_drop}
\end{equation}
For constant half-height $h_0$, this reduces to the plane Poiseuille
pressure drop $3\mu Q(x_b-x_a)/(2h_0^3)$. At fixed flux, the normalized
hydraulic resistance estimate is
\begin{equation}
 \frac{\Delta p}{\Delta p_0}
 =\frac{h_0^3}{x_b-x_a}\int_{x_a}^{x_b}\frac{dx}{h(x)^3}.
\end{equation}
This is a limiting analytical check, not a simulation result for the
proposed severity values. Its reliability requires sufficiently small wall
slope and inertial effects; a full Navier--Stokes solution is needed when
those assumptions fail \citep{batchelor1967}.

\section{Baseline trajectory criteria}\label{sec:trajectory_appendix}
\Cref{tab:path_limits} retains the seven baseline comparisons associated
with the field results in \cref{ch:basic}. RMS position discrepancy is
normalized by $W=0.3$. The current numerical limits require discrepancy
below 5\%, common-time coverage above 95\%, and agreement of every endpoint
class. Some individual rows satisfy those conditions, but the complete set
does not. The trajectories are consequently not used as validated
deposition evidence.

\begin{table}[htbp]\centering\small
\caption{Complete baseline particle comparison; no initial position is
removed to improve the aggregate outcome.}\label{tab:path_limits}
\begin{tabularx}{\textwidth}{@{}rrrY@{}}\toprule
Initial $\eta$ & RMS position/$W$ (\%) & Common time (\%) & Same endpoint class\\\midrule
\EvidenceRows{trajectories}
\bottomrule\end{tabularx}
\end{table}

\section{Planar mass convention}
For a physical sphere of diameter $d$, mass is $\rho_p\pi d^3/6$.
A parcel containing $N_i$ such spheres has mass
$N_i\rho_p\pi d^3/6$. Dividing by $\rho_fH^2B$ gives
\begin{equation}
 \frac{N_i\rho_p\pi d^3/6}{\rho_fH^2B}
 =\frac{N_iR\pi(d/H)^3}{6(B/H)},
\end{equation}
which is \cref{eq:parcel_mass}. The dimensional reaction source in a planar
fluid equation has force divided by depth and area. This is why a
two-dimensional Gaussian normalization uses area rather than volume. The
reference depth records the physical interpretation without introducing a
resolved spanwise flow field.

\chapter{Computational and Reproducibility Record}\label{app:reproducibility}

\section{Completed numerical records}
The baseline results are associated with the computational record
\path{basic_particle_problem/results/final_study}. The Gaussian
verification results and the diagnostic coupled calculations are associated
with \path{final_cloud_code/outputs/evaluated_2d_final}. These two result
sets are kept separate from short environment smoke tests. The smoke tests
are not used as research results in this thesis.

The archived repository revisions are
\begin{quote}\small\ttfamily
basic\_particle\_problem:\par
c3eb0efe76bcd62a3299dcd6c28922dbfa4c667a\par\medskip
final\_cloud\_code:\par
9d996c4f555bf5bf09e9f6fa78a3877e44310ec9
\end{quote}
The source bundle contains a SHA-256 manifest of the numerical inputs and
method files used to prepare the manuscript. The original result files are
copied without numerical modification. Figure scripts change presentation,
not the stored solution values. Gauge alignment and the exact computations
used to recompute metrics are explicit in those scripts.

The historical result manifests retain their own source hashes. The
archived baseline model and geometry, and the Gaussian model, physics, and
verification modules match their recorded hashes. Current driver and
configuration files do not all match the historical manifests; they are
retained for method inspection and are not claimed as an exact replay of
the original training execution.

\section{Training and evaluation distinction}
The baseline final evaluation restores previously trained network weights.
The associated original flow configuration uses 1,000 Adam epochs, a
250-iteration L-BFGS budget, and a refinement budget of 1,800 epochs with
700 L-BFGS iterations. The original force budget is 1,000 epochs. The two
flow seeds are 42 and 43. The later evaluation configuration contains
different default training fields, which do not imply that restored weights
were retrained with those defaults.

The Gaussian channel uses seed 7, float64, and the evaluated configuration
with 1,200 fluid/parcel Adam steps and a 200-iteration L-BFGS budget. Its
reduced ODE tests use 700 Adam steps with a 120-iteration refinement budget;
the Taylor--Green test uses 1,200 Adam steps with a 200-iteration refinement
budget. An iteration budget need not equal the number of closure
evaluations in a line-search optimizer. The stored status files and
configuration distinguish these quantities.

Reference values are excluded from the Gaussian reduced-test optimization.
In V01 the initial response is imposed. In V02 the initial values and total
momentum relation are imposed. In V10 the initial velocity and pressure
structure are imposed while the temporal decay is held out. These choices
determine which aspects of the tests are independent predictions.

\section{Software and execution scope}
The numerical environment is Linux with Python 3.12. The baseline CPU
environment uses PyTorch 2.14.1, and the Gaussian CPU environment uses
PyTorch 2.6.0. The scientific stack includes NumPy 2.3.5, SciPy 1.17.0,
Matplotlib 3.10.8, and scikit-fem 12.0.2. The preserved historical Gaussian
runtime file records CPU execution with two threads. These are
computational reproducibility details; no GPU performance or speedup is
claimed.

The source bundle supports two levels of reproduction. Compiling
\path{thesis.tex} uses the included tables and vector figures and does not
require retraining. Rebuilding the figures additionally uses the archived
arrays, model definitions, and two small baseline checkpoints. Independent
metric checks include the full Poiseuille velocity vector, centered
pressure, the final held-out drag evaluation, and the reduced Gaussian
benchmarks.

Reproducing all training runs requires the complete corresponding
repositories and their dependency environments. It also requires their
reference-input files. The manuscript bundle is a reproducible evidence
and typesetting package, not a claim that every full simulation can be
retrained from the manuscript alone.

\section{Figures and numerical transformations}
Nineteen figures are included, each registered with its input source or
its schematic status. Architecture figures are explanatory drawings.
The constriction and kernel-wall illustrations specify design choices.
They are never registered as computed flow results. Numerical plots retain
all evaluation arrays for aggregate errors; marker subsampling is used only
to make reference curves readable.

Pressure comparisons subtract an additive constant independently from each
array. Paired field panels use common color limits. Discrepancy panels have
their own explicitly described scale and are not smoothed to conceal local
error. The drag parity plot uses a stored diagnostic set different from
the final held-out metric set; both are identified in the caption and
validation record.

Five additional figures document the finite coupled parcel response of the
identified $\St=0.6$, $\beta=0.05$ calculation. Their diagnostic status and
the unresolved fluid residual are explicit in \cref{ch:gaussian}; they
are not relabeled as accepted predictive evidence. No synthetic
demonstration or extruded three-dimensional illustration is used as a
simulation result.
The material limitations are retained in the main chapters and the complete
baseline trajectory table. Additional unsuccessful run histories remain
outside the figure selection because they do not add accepted transport
results.


\end{document}
